% Chapter of the map: The structure of NP and dichotomies.
% Conventions and the template of an entry: blueprint/README.md. Check with
% blueprint/tools/check_names.sh. Labels of this file begin with kind:structure (def:structure-..., thm:structure-...).
\chapter{The structure of NP and dichotomies}
\label{chap:structure}

This chapter collects the statements of the map about what lies inside $\cNP$ and how it divides.
It starts with problems of intermediate difficulty (Ladner's construction, the clock it needs, and the
statement that an intermediate problem exists if and only if $\cP \neq \cNP$), goes on to the
algebraic description of Boolean constraint problems by polymorphisms and the dichotomy it gives
under named assumptions, and then lists concrete problems placed in $\cNP$. After that come the
pumping move against $\SAT$, lower bounds in the model of a search tree and of cutting planes, and
statements about proof systems and about the rate of approximation. The last four sections are about
padding, the three switches, and the diagonal argument over enumerations. Each section says in which
model its statements live: Turing machines of cslib-frontier, the abstract classes of the map, finite
structures, or an abstract cost model.

\section{Problems between $\cP$ and the $\cNP$-complete ones}
\label{sec:structure-ladner}

The statements of this section are about decision problems over binary words, with the classes
$\cP$ and $\cNP$ and the notion of $\cNP$-completeness of cslib-frontier (Part~I). The first
group is Ladner's theorem in three forms; the second group is the clock that the construction
needs.

\begin{definition}[Intermediate language]
  \label{def:structure-intermedio}
  \lean{Sagitta.Map.Borde.FractalSAT.Intermedio}
  \leanok
  \uses{def:P, def:NP, def:reduction}
  A set $L$ of binary words is \emph{intermediate} when $L \in \cNP$ (over binary words), $L
  \notin \cP$ and $L$ is not $\cNP$-complete.
  \own{}
  \leanfile{Sagitta/Map/Borde/FractalSAT.lean}
\end{definition}

\begin{theorem}[Ladner's theorem, as stated over binary words in the map]
  \label{thm:structure-ladner-cslib}
  \lean{Sagitta.Map.Borde.RelojDeLadner.ladner_cslib}
  \leanok
  \uses{thm:ladner}
  Let $\alpha$ be a type and $\enc : \alpha \hookrightarrow \mathrm{List}\,\mathrm{Bool}$ an
  injective encoding. If $\cP \neq \cNP$ for $\enc$, then there is a set $L$ of binary words with
  $L \in \cNP$, $L \notin \cP$ and $L$ not $\cNP$-complete.
  \source{\cite{Ladner75}.}
  \fidelity{Not checked against the printed source.}
  \leanfile{Sagitta/Map/Borde/RelojDeLadner.lean}
\end{theorem}

\begin{proof}
  \leanok
  It is the theorem of cslib-frontier (Theorem~\ref{thm:ladner}) applied as it stands.
\end{proof}

\begin{theorem}[Ladner's language outside a set in $\cP$, by words]
  \label{thm:structure-ladner-palabras}
  \lean{Sagitta.Map.Borde.BordeAlto.ladner_por_palabras_SATH}
  \leanok
  \uses{thm:structure-ladner-cslib, thm:ladner, def:P, def:NP, def:reduction}
  Let $\enc : \alpha \hookrightarrow \mathrm{List}\,\mathrm{Bool}$ be an encoding with $\cP \neq
  \cNP$ for $\enc$. Then there is a set $G$ of binary words with $G \in \cP$ such that the set
  of words of $\mathrm{SATH}$ that are not in $G$ is in $\cNP$, is not in $\cP$ and is not
  $\cNP$-complete. Here $\mathrm{SATH}$ is the language of cslib-frontier used in its proof of
  Ladner's theorem.
  \source{\cite{Ladner75}.}
  \fidelity{Not checked against the printed source.}
  \leanfile{Sagitta/Map/Borde/BordeAlto.lean}
\end{theorem}

\begin{proof}
  \leanok
  It applies a lemma of the same file (for every $B \in \cNP$ outside $\cP$ there is such a $G$,
  from cslib-frontier) to the facts that $\mathrm{SATH}$ is in $\cNP$ and, if $\cP \neq \cNP$,
  not in $\cP$.
\end{proof}

\begin{theorem}[An intermediate language exists if $\cP \neq \cNP$]
  \label{thm:structure-ladner-fractal}
  \lean{Sagitta.Map.Borde.FractalSAT.ladner}
  \leanok
  \uses{def:structure-intermedio, thm:structure-ladner-cslib, def:P, def:NP}
  If $\cP \neq \cNP$ over binary words (the classes of cslib-frontier for the binary encoding of
  words), then some language is intermediate.
  \source{\cite{Ladner75}.}
  \fidelity{Not checked against the printed source.}
  \leanfile{Sagitta/Map/Borde/FractalSAT.lean}
\end{theorem}

\begin{proof}
  \leanok
  Theorem~\ref{thm:structure-ladner-cslib} with the binary encoding of words.
\end{proof}

\subsection{An intermediate language exists if and only if $\cP \neq \cNP$}

\begin{theorem}[Intermediate languages exist if and only if $\cP \neq \cNP$]
  \label{thm:structure-intermedio-iff}
  \lean{Sagitta.Map.Borde.FractalSAT.intermedio_iff_P_ne_NP}
  \leanok
  \uses{def:structure-intermedio, thm:structure-ladner-fractal}
  There is an intermediate language if and only if $\cP \neq \cNP$ over binary words.
  \own{}
  \leanfile{Sagitta/Map/Borde/FractalSAT.lean}
\end{theorem}

\begin{proof}
  \leanok
  An intermediate language is in $\cNP$ and not in $\cP$, so the two classes differ; the converse
  is Theorem~\ref{thm:structure-ladner-fractal}.
\end{proof}

\begin{theorem}[If $\cP = \cNP$, every non-trivial language of $\cNP$ is complete]
  \label{thm:structure-completo-de-p-eq-np}
  \lean{Sagitta.Map.Borde.FractalSAT.completo_de_P_eq_NP}
  \leanok
  \uses{def:P, def:NP, def:reduction}
  Assume that $\cP = \cNP$ for every type $\alpha$ and every encoding $\enc : \alpha
  \hookrightarrow \mathrm{List}\,\mathrm{Bool}$. Let $\beta$ be a type with an encoding $\enc_B$, let
  $L \in \cNP$ for $\enc_B$ be a set of elements of $\beta$, and let $y_1, y_0 \in \beta$ with
  $y_1 \in L$ and $y_0 \notin L$. Then $L$ is $\cNP$-complete for $\enc_B$.
  \own{}
  \leanfile{Sagitta/Map/Borde/FractalSAT.lean}
\end{theorem}

\begin{proof}
  \leanok
  A problem $A \in \cNP$ is in $\cP$, so a function computable in polynomial time decides it; the
  reduction sends $a$ to $y_1$ if that function accepts and to $y_0$ otherwise.
\end{proof}

\subsection{The clock of Ladner's construction}

\begin{definition}[Eventually homogeneous set of lengths]
  \label{def:structure-homogeneo}
  \lean{Sagitta.Map.Borde.RelojDeLadner.HomogeneoEnLargo}
  \leanok
  A set $B \subseteq \mathbb{N}$ is \emph{eventually homogeneous} when either $n \in B$ for all
  large $n$, or $n \notin B$ for all large $n$.
  \own{}
  \leanfile{Sagitta/Map/Borde/RelojDeLadner.lean}
\end{definition}

\begin{proposition}[A set of lengths is not eventually homogeneous if and only if it and its complement are infinite]
  \label{prop:structure-no-homogeneo}
  \lean{Sagitta.Map.Borde.RelojDeLadner.no_homogeneo_iff}
  \leanok
  \uses{def:structure-homogeneo}
  For $B \subseteq \mathbb{N}$, $B$ is not eventually homogeneous if and only if $B$ is infinite and
  its complement $\mathbb{N} \setminus B$ is infinite.
  \own{}
  \leanfile{Sagitta/Map/Borde/RelojDeLadner.lean}
\end{proposition}

\begin{proof}
  \leanok
  ``For all large $n$, $n \in B$'' says that the complement is finite, and ``for all large $n$,
  $n \notin B$'' says that $B$ is finite; negate the disjunction.
\end{proof}

\begin{definition}[Cheap clock]
  \label{def:structure-reloj-barato}
  \lean{Sagitta.Map.Borde.RelojDeLadner.RelojBarato}
  \leanok
  \uses{def:polytime}
  A function $e : \mathbb{N} \to \mathbb{N}$ is a \emph{cheap clock} when the function that sends
  a pair $(m, n)$ to the truth value of $e(m) \le n$ is computable in polynomial time, with the
  pair and both numbers written in unary.
  \own{}
  \leanfile{Sagitta/Map/Borde/RelojDeLadner.lean}
\end{definition}

\begin{proposition}[The doubling clock is cheap]
  \label{prop:structure-reloj-dobles}
  \lean{Sagitta.Map.Borde.RelojDeLadner.relojBarato_dobles}
  \leanok
  \uses{def:structure-reloj-barato}
  The function $m \mapsto 2^m - 1$ is a cheap clock.
  \own{}
  \leanfile{Sagitta/Map/Borde/RelojDeLadner.lean}
\end{proposition}

\begin{proof}
  \leanok
  The test $2^m - 1 \le n$ equals a comparison of a capped power of two with its cap, which is
  computable in polynomial time by the tactic of cslib-frontier; the equality of the two
  functions is a computation with natural numbers.
\end{proof}

\begin{theorem}[Ladner's clock is computable in polynomial time, in unary]
  \label{thm:structure-reloj-de-ladner}
  \lean{Sagitta.Map.Borde.RelojDeLadner.reloj_de_ladner}
  \leanok
  \uses{def:polytime}
  The function $H : \mathbb{N} \to \mathbb{N}$ of cslib-frontier's proof of Ladner's theorem is
  computable in polynomial time, with input and output in unary.
  The function $H$ is the one of Ladner's construction, as formalized in cslib-frontier.
  \own{}
  \leanfile{Sagitta/Map/Borde/RelojDeLadner.lean}
\end{theorem}

\begin{proof}
  \leanok
  It is the theorem of cslib-frontier that $H$ is computable in polynomial time.
\end{proof}

\section{Polymorphisms of Boolean constraint languages}
\label{sec:structure-polymorphisms}

The statements of this section are about finite Boolean relations and the operations on
$\{0,1\}$ that preserve them. There is no machine and no cost measure in them: they describe
the algebra that decides, in the next section, which constraint problems are in $\cP$.

\begin{definition}[Boolean relations, preservation and polymorphisms]
  \label{def:structure-polymorphism}
  \lean{Sagitta.Map.Informacion.Polimorfismos.Relacion, Sagitta.Map.Informacion.Polimorfismos.Relaciones, Sagitta.Map.Informacion.Polimorfismos.Preserva, Sagitta.Map.Informacion.Polimorfismos.Polimorfismo, Sagitta.Map.Informacion.Polimorfismos.proy}
  \leanok
  A Boolean relation of arity $r$ is a predicate on $\{0,1\}^r$ (functions $\mathrm{Fin}\,r \to
  \mathrm{Bool}$); a \emph{language} $\Gamma$ is a set of Boolean relations of arbitrary arities.
  An operation $f : \{0,1\}^k \to \{0,1\}$ \emph{preserves} an $r$-ary relation $R$ when, for
  every choice of $k$ tuples $t_1, \dots, t_k$ of $R$, the tuple whose coordinate $c$ is
  $f(t_1(c), \dots, t_k(c))$ is in $R$. It is a \emph{polymorphism} of $\Gamma$ when it preserves
  every relation of $\Gamma$. The projection $\mathrm{proy}_j$ sends $x$ to $x_j$.
  \own{}
  \leanfile{Sagitta/Map/Informacion/Polimorfismos.lean}
\end{definition}

\begin{definition}[Clauses, 2-SAT, 3-SAT, Horn and affine languages; majority, minority and AND]
  \label{def:structure-languages}
  \lean{Sagitta.Map.Informacion.Polimorfismos.clausula, Sagitta.Map.Informacion.Polimorfismos.dosSAT, Sagitta.Map.Informacion.Polimorfismos.tresSAT, Sagitta.Map.Informacion.Polimorfismos.paridad, Sagitta.Map.Informacion.Polimorfismos.ecuacion, Sagitta.Map.Informacion.Polimorfismos.xorSAT, Sagitta.Map.Informacion.Polimorfismos.maj, Sagitta.Map.Informacion.Polimorfismos.maltsev, Sagitta.Map.Informacion.Polimorfismos.y}
  \leanok
  \uses{def:structure-polymorphism}
  For $\sigma : \mathrm{Fin}\,r \to \mathrm{Bool}$, the clause $\mathrm{clausula}\,\sigma$ is the
  relation that holds of $x$ when $x_c = \sigma_c$ for some coordinate $c$. The language
  $\mathrm{dosSAT}$ is the set of clauses with $r = 2$, $\mathrm{tresSAT}$ the set of clauses with
  $r = 3$. The parity of a tuple is the exclusive or of its coordinates; $\mathrm{ecuacion}\,r\,b$
  is the relation ``the parity of $x$ is $b$'', and $\mathrm{xorSAT}$ is the set of these
  relations with $r \in \{2, 3\}$ and $b \in \{0,1\}$. The operations are the majority
  $\mathrm{maj}(x_0,x_1,x_2) = (x_0 \wedge x_1) \vee (x_0 \wedge x_2) \vee (x_1 \wedge x_2)$, the
  minority $\mathrm{maltsev}(x_0,x_1,x_2) = x_0 \oplus x_1 \oplus x_2$, and the conjunction
  $y(x_0,x_1) = x_0 \wedge x_1$.
  \own{}
  \leanfile{Sagitta/Map/Informacion/Polimorfismos.lean}
\end{definition}

\begin{theorem}[The polymorphisms of 3-SAT of arity at most three are the projections]
  \label{thm:structure-tressat-pol}
  \lean{Sagitta.Map.Informacion.Polimorfismos.tresSAT_pol_iff}
  \leanok
  \uses{def:structure-polymorphism, def:structure-languages}
  Let $k$ be a natural number with $0 < k \le 3$ and $f : \{0,1\}^k \to \{0,1\}$. Then $f$ is a
  polymorphism of $\mathrm{tresSAT}$ if and only if $f$ is a projection $\mathrm{proy}_j$ for some
  $j$.
  \source{\cite{Schaefer78,Jeavons98}.}
  \fidelity{Not checked against the printed source.}
  \leanfile{Sagitta/Map/Informacion/Polimorfismos.lean}
\end{theorem}

\begin{proof}
  \leanok
  If $f$ preserves the eight clauses, then on any choice of three inputs it agrees with a
  projection; taking the unit inputs gives a coordinate $j$ with $f(e_j) = 1$, and that forces $f
  = \mathrm{proy}_j$. Conversely every projection is a polymorphism of every language.
\end{proof}

\begin{theorem}[Majority is a polymorphism of 2-SAT]
  \label{thm:structure-maj-dossat}
  \lean{Sagitta.Map.Informacion.Polimorfismos.maj_polimorfismo_dosSAT}
  \leanok
  \uses{def:structure-polymorphism, def:structure-languages}
  The majority operation is a polymorphism of $\mathrm{dosSAT}$.
  \source{\cite{Schaefer78,Jeavons98}.}
  \fidelity{Not checked against the printed source.}
  \leanfile{Sagitta/Map/Informacion/Polimorfismos.lean}
\end{theorem}

\begin{proof}
  \leanok
  Among three pairs of Booleans, the majority taken coordinate by coordinate is one of the three
  pairs, so majority preserves every binary relation; this is checked by \texttt{decide}.
\end{proof}

\begin{theorem}[The minority operation is a polymorphism of the affine language]
  \label{thm:structure-maltsev-xorsat}
  \lean{Sagitta.Map.Informacion.Polimorfismos.maltsev_polimorfismo_xorSAT}
  \leanok
  \uses{def:structure-polymorphism, def:structure-languages}
  The operation $x \oplus y \oplus z$ is a polymorphism of $\mathrm{xorSAT}$.
  \source{\cite{Schaefer78,Jeavons98}.}
  \fidelity{Not checked against the printed source.}
  \leanfile{Sagitta/Map/Informacion/Polimorfismos.lean}
\end{theorem}

\begin{proof}
  \leanok
  Parity commutes with the exclusive or of three tuples coordinate by coordinate, and $b \oplus b
  \oplus b = b$, so the operation preserves each equation.
\end{proof}

\section{The Boolean dichotomy, under named assumptions}
\label{sec:structure-dichotomy}

The statements of this section are about constraint problems $\mathrm{CSP}(\Gamma)$ for a finite
Boolean language $\Gamma$. An instance is a list of constraints, each a relation of $\Gamma$ with a
tuple of variables. The predicates ``$\mathrm{CSP}(\Gamma)$ is in $\cP$'' and ``is
$\cNP$-complete'' are not defined here: they are arbitrary predicates $\mathrm{EnP}$ and
$\mathrm{NPc}$ on languages, and what is assumed about them is gathered in one hypothesis.

\begin{theorem}[The polymorphisms of 3-SAT are the projections, in every arity]
  \label{thm:structure-tressat-pol-todas}
  \lean{Sagitta.Map.Borde.DicotomiaBooleana.tresSAT_pol_iff_todas}
  \leanok
  \uses{thm:structure-tressat-pol}
  Let $k$ be a natural number and $f : \{0,1\}^k \to \{0,1\}$. Then $f$ is a polymorphism of
  $\mathrm{tresSAT}$ if and only if $f = \mathrm{proy}_j$ for some $j$.
  \own{}
  \leanfile{Sagitta/Map/Borde/DicotomiaBooleana.lean}
\end{theorem}

\begin{proof}
  \leanok
  The sets $A$ of coordinates with $f(\mathbf{1}_A) = 1$ contain the full set, not the empty set,
  and for each $A$ in it and $j \in A$ contain $\{j\}$ or $A \setminus \{j\}$. Induction on the
  size of $A$ gives a coordinate $j$ with $f(e_j) = 1$, and then $f = \mathrm{proy}_j$.
\end{proof}

\begin{definition}[Primitive positive definitions and the canonical one]
  \label{def:structure-pp}
  \lean{Sagitta.Map.Borde.DicotomiaBooleana.DefPP, Sagitta.Map.Borde.DicotomiaBooleana.DefPP.Define, Sagitta.Map.Borde.DicotomiaBooleana.canonica}
  \leanok
  \uses{def:structure-polymorphism}
  A primitive positive definition of an $s$-ary relation over a language $\Gamma$ consists of a
  finite type $W$ of variables, a finite family of relations of $\Gamma$ with, for each, a tuple of
  variables of $W$, and $s$ distinguished variables of $W$ (which may repeat). It defines the
  relation of those $y \in \{0,1\}^s$ for which some assignment $x : W \to \{0,1\}$ satisfies all
  the relations of the family and gives the distinguished variables the values of $y$. For a finite
  $\Gamma$ and an $s$-ary relation $S$, the canonical definition has as variables the functions
  $\{0,1\}^{m}$ with $m$ the number of tuples of $S$, as constraints all those of $\Gamma$ that hold
  at the $m$ canonical assignments, and as $c$-th distinguished variable the column $c$ of the
  list of tuples of $S$.
  \own{}
  \leanfile{Sagitta/Map/Borde/DicotomiaBooleana.lean}
\end{definition}

\begin{theorem}[Galois connection for finite Boolean languages]
  \label{thm:structure-galois}
  \lean{Sagitta.Map.Borde.DicotomiaBooleana.galois}
  \leanok
  \uses{def:structure-pp, def:structure-polymorphism}
  Let $\Gamma$ be a finite language and $S$ an $s$-ary Boolean relation. Assume that every
  polymorphism of $\Gamma$, of every arity $k$, preserves $S$. Then for every $y \in \{0,1\}^s$,
  $S(y)$ holds if and only if the canonical definition of $S$ over $\Gamma$ defines $y$.
  \source{\cite{Geiger68,BodnarchukEtAl69}.}
  \fidelity{Not checked against the printed source.}
  \leanfile{Sagitta/Map/Borde/DicotomiaBooleana.lean}
\end{theorem}

\begin{proof}
  \leanok
  A tuple of $S$ is defined by the assignment that sends each canonical variable to its $i$-th
  coordinate. Conversely a defining assignment is a polymorphism of $\Gamma$ of arity $m$; by
  hypothesis it preserves $S$ and is applied to the list of tuples of $S$.
\end{proof}

\begin{theorem}[Replacing each constraint by its definition preserves satisfiability]
  \label{thm:structure-traducir}
  \lean{Sagitta.Map.Borde.DicotomiaBooleana.traducir_correcto}
  \leanok
  \uses{def:structure-pp}
  Let $\Gamma, \Delta$ be languages and $V$ a type of variables. Suppose that every relation $R
  \in \Delta$ has a primitive positive definition over $\Gamma$ whose distinguished variables are
  pairwise distinct (injective) and which defines $R$. Then an instance $I$ of
  $\mathrm{CSP}(\Delta)$ over $V$ is satisfiable if and only if its translation, in which each
  constraint is replaced by the constraints of its definition with fresh existential variables, is
  satisfiable as an instance of $\mathrm{CSP}(\Gamma)$.
  \own{}
  \leanfile{Sagitta/Map/Borde/DicotomiaBooleana.lean}
\end{theorem}

\begin{proof}
  \leanok
  From a solution of $I$ take, for each constraint, the assignment that witnesses its definition,
  and join them; from a solution of the translation restrict to the original variables and read
  the same witnesses back.
\end{proof}

\begin{definition}[The six identities of Schaefer]
  \label{def:structure-six}
  \lean{Sagitta.Map.Borde.DicotomiaBooleana.TieneUnaDeLasSeis, Sagitta.Map.Borde.DicotomiaBooleana.const, Sagitta.Map.Borde.DicotomiaBooleana.o}
  \leanok
  \uses{def:structure-polymorphism, def:structure-languages}
  A language $\Gamma$ \emph{has one of the six} when at least one of the following is a
  polymorphism of $\Gamma$: the constant $0$, the constant $1$, the conjunction $x_0 \wedge x_1$,
  the disjunction $x_0 \vee x_1$, the majority, the minority $x_0 \oplus x_1 \oplus x_2$.
  \own{}
  \leanfile{Sagitta/Map/Borde/DicotomiaBooleana.lean}
\end{definition}

\begin{definition}[Hypothesis: the assumptions of the Boolean dichotomy]
  \label{def:structure-supuestos}
  \lean{Sagitta.Map.Borde.DicotomiaBooleana.Supuestos}
  \leanok
  \uses{def:structure-six, def:structure-pp}
  For predicates $\mathrm{EnP}$ and $\mathrm{NPc}$ on languages, $\mathrm{Supuestos}$ is the
  conjunction of the following six statements.
  \begin{itemize}
    \item (easy side) For every finite $\Gamma$ that has one of the six, $\mathrm{EnP}(\Gamma)$.
    \item (Post) For every $\Gamma$ that has none of the six, every polymorphism $f$ of $\Gamma$ of
      any arity is essentially unary: $f = \mathrm{proy}_j$ or $f(x) = \lnot x_j$ for some $j$.
    \item $\mathrm{NPc}(\mathrm{tresSAT})$.
    \item $\mathrm{NPc}$ of the language that contains only the relation ``not all equal'' on
      three variables.
    \item (reduction) For finite $\Gamma$ and $\Delta$ with $\mathrm{NPc}(\Delta)$, if every
      relation of $\Delta$ has an injective-variable primitive positive definition over $\Gamma$
      that defines it, then $\mathrm{NPc}(\Gamma)$.
    \item (separation) For every $\Gamma$, $\mathrm{NPc}(\Gamma)$ implies not $\mathrm{EnP}(\Gamma)$.
  \end{itemize}
  The hypothesis is a structure of type \texttt{Prop}; this entry says that it is defined, not that
  it holds.
  \own{}
  \leanfile{Sagitta/Map/Borde/DicotomiaBooleana.lean}
\end{definition}

\begin{theorem}[Dichotomy for finite Boolean languages, under the assumptions]
  \label{thm:structure-dicotomia}
  \lean{Sagitta.Map.Borde.DicotomiaBooleana.dicotomia}
  \leanok
  \uses{def:structure-supuestos, def:structure-six}
  Assume $\mathrm{Supuestos}(\mathrm{EnP}, \mathrm{NPc})$. Then for every finite language $\Gamma$,
  either $\mathrm{EnP}(\Gamma)$ and not $\mathrm{NPc}(\Gamma)$, or $\mathrm{NPc}(\Gamma)$ and not
  $\mathrm{EnP}(\Gamma)$.
  \source{\cite{Schaefer78}.}
  \fidelity{Not checked against the printed source.}
  \leanfile{Sagitta/Map/Borde/DicotomiaBooleana.lean}
\end{theorem}

\begin{proof}
  \leanok
  If $\Gamma$ has one of the six, the easy side gives $\mathrm{EnP}$ and the separation excludes
  $\mathrm{NPc}$. Otherwise it is $\mathrm{NPc}$: when negation is a polymorphism the relation ``not
  all equal'' is defined over $\Gamma$ by its canonical definition, and when it is not, each clause
  of 3-SAT is; the reduction assumption closes both cases.
\end{proof}

\begin{theorem}[Membership in $\cP$ is having one of the six, under the assumptions]
  \label{thm:structure-interruptor}
  \lean{Sagitta.Map.Borde.DicotomiaBooleana.interruptor}
  \leanok
  \uses{def:structure-supuestos, def:structure-six}
  Assume $\mathrm{Supuestos}(\mathrm{EnP}, \mathrm{NPc})$. For every finite language $\Gamma$,
  $\mathrm{EnP}(\Gamma)$ if and only if $\Gamma$ has one of the six.
  \own{}
  \leanfile{Sagitta/Map/Borde/DicotomiaBooleana.lean}
\end{theorem}

\begin{proof}
  \leanok
  The easy side is the assumption; for the converse, if $\Gamma$ had none of the six it would be
  $\mathrm{NPc}$ (the case analysis above), which the separation assumption makes incompatible with
  $\mathrm{EnP}$.
\end{proof}

\begin{theorem}[No finite Boolean language is left in between, under the assumptions]
  \label{thm:structure-sin-intermedios}
  \lean{Sagitta.Map.Borde.DicotomiaBooleana.sin_intermedios}
  \leanok
  \uses{def:structure-supuestos, thm:structure-dicotomia}
  Assume $\mathrm{Supuestos}(\mathrm{EnP}, \mathrm{NPc})$. For every finite language $\Gamma$, it is
  not the case that both $\mathrm{EnP}(\Gamma)$ fails and $\mathrm{NPc}(\Gamma)$ fails.
  \own{}
  \leanfile{Sagitta/Map/Borde/DicotomiaBooleana.lean}
\end{theorem}

\begin{proof}
  \leanok
  By Theorem~\ref{thm:structure-dicotomia} one of the two holds.
\end{proof}

\subsection{Horn and dual-Horn satisfiability by unit propagation}

A Horn clause over $n$ variables is a body (a finite set of variables) and a head, which is a
variable or absent; a set $S$ of variables (those set to $1$) satisfies it when the body is not
contained in $S$, or the head is present and in $S$. A model of a list of clauses is a set that
satisfies all of them. The algorithm $\mathrm{decideHorn}$ propagates the heads of the clauses whose
body already holds, $n + 1$ times starting from the empty set, and checks the clauses on the set
that results. The statements are about these finite objects and carry no machine model.

\begin{definition}[Horn clauses and unit propagation]
  \label{def:structure-horn}
  \lean{Sagitta.Map.Borde.DicotomiaBooleana.ClausulaHorn, Sagitta.Map.Borde.DicotomiaBooleana.ModeloHorn, Sagitta.Map.Borde.DicotomiaBooleana.ModeloDual, Sagitta.Map.Borde.DicotomiaBooleana.decideHorn}
  \leanok
  A Horn clause over $n$ variables is a pair (body, head) with body a finite set of $\mathrm{Fin}\,n$
  and head an optional element of $\mathrm{Fin}\,n$. The model of a list $cs$ is a set $S$ that
  satisfies every clause as described above. A dual model of $cs$ is a set $S$ whose complement is
  a model. The algorithm $\mathrm{decideHorn}(cs)$ accepts if and only if every clause holds on the
  set obtained by $n + 1$ propagation steps from the empty set.
  \own{}
  \leanfile{Sagitta/Map/Borde/DicotomiaBooleana.lean}
\end{definition}

\begin{theorem}[Unit propagation decides Horn satisfiability]
  \label{thm:structure-horn-correcto}
  \lean{Sagitta.Map.Borde.DicotomiaBooleana.decideHorn_correcto}
  \leanok
  \uses{def:structure-horn}
  For every list $cs$ of Horn clauses over $n$ variables, $\mathrm{decideHorn}(cs)$ is true if and
  only if $cs$ has a model.
  \source{\cite{DowlingGallier84}.}
  \fidelity{Not checked against the printed source.}
  \leanfile{Sagitta/Map/Borde/DicotomiaBooleana.lean}
\end{theorem}

\begin{proof}
  \leanok
  Each propagation step that does not stay fixed adds a variable, so the set after $n + 1$ steps is
  a fixed point; it is contained in every model, so it is the least model, and the check on it is
  exactly the existence of a model.
\end{proof}

\begin{theorem}[The same algorithm decides dual-Horn satisfiability]
  \label{thm:structure-dual-correcto}
  \lean{Sagitta.Map.Borde.DicotomiaBooleana.decideDual_correcto}
  \leanok
  \uses{def:structure-horn, thm:structure-horn-correcto}
  For every list $cs$ of Horn clauses over $n$ variables, $\mathrm{decideHorn}(cs)$ is true if and
  only if $cs$ has a dual model.
  \source{\cite{DowlingGallier84}.}
  \fidelity{Not checked against the printed source.}
  \leanfile{Sagitta/Map/Borde/DicotomiaBooleana.lean}
\end{theorem}

\begin{proof}
  \leanok
  Taking complements is a bijection between models and dual models, so the statement is
  Theorem~\ref{thm:structure-horn-correcto}.
\end{proof}

\begin{proposition}[The models of a Horn list are closed under intersection]
  \label{prop:structure-modelo-inter}
  \lean{Sagitta.Map.Borde.DicotomiaBooleana.modelo_inter}
  \leanok
  \uses{def:structure-horn}
  If $S$ and $T$ are models of a list $cs$ of Horn clauses, then $S \cap T$ is a model of $cs$.
  \source{\cite{DowlingGallier84}.}
  \fidelity{Not checked against the printed source.}
  \leanfile{Sagitta/Map/Borde/DicotomiaBooleana.lean}
\end{proposition}

\begin{proof}
  \leanok
  Case analysis on whether the body of a clause is contained in $S \cap T$: if so it is contained
  in both, and the head, which is present, lies in both.
\end{proof}

\section{Concrete problems in $\cNP$}
\label{sec:structure-concrete}

The statements of this section place three concrete problems in $\cNP$: shortest common
superstring, primality and the graph of the least prime factor, and satisfiability restricted to
a set of lengths. The first two live in the class $\cNP$ of the map, whose problems are functions
from binary words to Booleans and whose verifiers are computed by Turing machines of the Mathlib
kind in polynomial time; the third lives in the class $\cNP$ of cslib-frontier for the encoding of
CNF formulas.

\begin{theorem}[Shortest common superstring is in $\cNP$, if its verifier is computable in polynomial time]
  \label{thm:structure-supercadena}
  \lean{Sagitta.Map.Supercadena.en_NP}
  \leanok
  % uses-lean: Sagitta.Map.NP, Sagitta.Map.EsPolinomial
  Let the instance be read from a binary word as a bound $k$ (a first self-delimited block, in
  binary) followed by a list of binary words (the remaining blocks). The problem asks whether some
  binary word of length at most $k$ contains every word of the list as a contiguous factor. The
  verifier takes the input and a certificate $y$ and accepts if and only if $|y| \le k$ and every
  word of the list is a factor of $y$. Assume that the verifier is computable in polynomial time,
  that is, that a Turing machine of the Mathlib kind computes it in polynomial time. Then the problem
  is in the class $\cNP$ of the map.
  \source{\cite{GallantMaierStorer80}.}
  \fidelity{Not checked against the printed source.}
  \leanfile{Sagitta/Map/Supercadena.lean}
\end{theorem}

\begin{proof}
  \leanok
  The verifier decides the problem with a certificate no longer than the input: if a superstring
  of length at most $k$ exists, either the given one is already short or the concatenation of all
  the words, whose length is at most the length of the input, is also a superstring of length at
  most $k$. The polynomial bound of the certificate is $X$.
\end{proof}

\begin{definition}[Pratt certificates and the check of the graph of the least prime factor]
  \label{def:structure-pratt}
  \lean{Sagitta.Map.Acotada.Pratt.entradaOK, Sagitta.Map.Acotada.Pratt.prattOK, Sagitta.Map.Acotada.Pratt.revisaGrafo, Sagitta.Map.Acotada.Pratt.grafoMinFac, Sagitta.Map.Acotada.Pratt.primos}
  \leanok
  A flat Pratt certificate is a list of entries $(p, a, [q_1, \dots, q_k])$. An entry is correct
  inside a certificate $C$ when $p \ge 2$, $\prod q_i = p - 1$, $a^{p-1} \equiv 1 \pmod p$, and each
  $q_i$ is the first component of some entry of $C$ and satisfies $a^{(p-1)/q_i} \not\equiv 1 \pmod
  p$; the certificate is correct when all its entries are. The check $\mathrm{revisaGrafo}\,n\,m$ of
  a pair $(L, C)$ accepts: for $n = 0$ when $m = 2$, for $n = 1$ when $m = 1$, and otherwise when
  $\prod L = n$, $m \in L$, every $q \in L$ satisfies $m \le q$ and is the first component of an
  entry of $C$, and $C$ is correct. The problem $\mathrm{primos}$ asks whether the number coded by
  a binary word $w$ (the number $\mathrm{numeroDe}(w) - 1$) is prime; $\mathrm{grafoMinFac}$ asks
  whether, writing $\mathrm{numeroDe}(w) - 1$ as the pair $\langle n, m \rangle$ of Mathlib, $m$ is the
  least prime factor of $n$.
  \own{}
  \leanfile{Sagitta/Map/Borde/Pratt.lean}
\end{definition}

\begin{theorem}[Pratt certificates for the least prime factor have quadratic length]
  \label{thm:structure-certificado-grafo}
  \lean{Sagitta.Map.Acotada.Pratt.certificado_grafo}
  \leanok
  \uses{def:structure-pratt}
  For every natural number $n$ there is a pair $c = (L, C)$ such that
  $\mathrm{revisaGrafo}\,n\,(\mathrm{minFac}\,n)\,c$ accepts, and the binary encoding of $c$ has
  length at most $63\,(\mathrm{size}(n) + 1)^2$.
  \source{\cite{Pratt75}.}
  \fidelity{Not checked against the printed source.}
  \leanfile{Sagitta/Map/Borde/Pratt.lean}
\end{theorem}

\begin{proof}
  \leanok
  Take $L$ the list of prime factors of $n$ and, from the existence of a certificate for each
  prime with at most $2\log_2 p - 1$ entries, a certificate $C$ for all of them; the sum of the
  sizes of the factors is bounded by the size of $n$ plus the number of factors, which gives the
  quadratic bound. For $n < 2$ the pair $([], [])$ works.
\end{proof}

\begin{theorem}[Every prime has a short certificate: the primes are in $\cNP$]
  \label{thm:structure-primos-np}
  \lean{Sagitta.Map.Acotada.Pratt.primos_NP}
  \leanok
  \uses{def:structure-pratt}
  % uses-lean: Sagitta.Map.NP
  The problem $\mathrm{primos}$ is in the class $\cNP$ of the map.
  \source{\cite{Pratt75}.}
  \fidelity{Not checked against the printed source.}
  \leanfile{Sagitta/Map/Borde/Pratt.lean}
\end{theorem}

\begin{proof}
  \leanok
  A prime has a correct certificate by the converse of Lucas's criterion (the units modulo a
  prime form a cyclic group) with at most $2\log_2 p - 1$ entries of $O(|p|)$ bits each, hence of
  quadratic length; checking a certificate is computable in polynomial time, and a correct
  certificate implies primality by Lucas's criterion, by strong induction on $p$.
\end{proof}

\begin{theorem}[The graph of the least prime factor is in $\cNP$]
  \label{thm:structure-grafo-minfac-np}
  \lean{Sagitta.Map.Acotada.Pratt.grafoMinFac_NP}
  \leanok
  \uses{def:structure-pratt, thm:structure-certificado-grafo}
  % uses-lean: Sagitta.Map.NP
  The problem $\mathrm{grafoMinFac}$ is in the class $\cNP$ of the map.
  \own{}
  \leanfile{Sagitta/Map/Borde/Pratt.lean}
\end{theorem}

\begin{proof}
  \leanok
  The verifier decodes the pair and the certificate and runs $\mathrm{revisaGrafo}$, which is
  computable in polynomial time; the bound $63(|w|+2)^2$ on the certificate comes from
  Theorem~\ref{thm:structure-certificado-grafo}.
\end{proof}

\begin{definition}[Satisfiability restricted to a set of lengths]
  \label{def:structure-sat-en}
  \lean{Sagitta.Map.Borde.Goodstein.satEn}
  \leanok
  \uses{def:threesat}
  For a test $c : \mathbb{N} \to \mathrm{Bool}$, $\mathrm{satEn}\,c$ is the set of CNF formulas
  $\varphi$ over $\mathbb{N}$ such that $c$ accepts the length of the encoding of $\varphi$ and
  $\varphi \in \SAT$.
  \own{}
  \leanfile{Sagitta/Map/Borde/Goodstein.lean}
\end{definition}

\begin{definition}[Hypothesis: the test of lengths is computable in polynomial time]
  \label{def:structure-longitudes-polinomiales}
  \lean{Sagitta.Map.Borde.Goodstein.LongitudesPolinomiales}
  \leanok
  \uses{def:polytime}
  For a test $c : \mathbb{N} \to \mathrm{Bool}$, $\mathrm{LongitudesPolinomiales}\,c$ says that the
  function that sends a CNF formula $\varphi$ to $c$ applied to the length of its encoding is
  computable in polynomial time, with the encoding of CNF formulas of cslib-frontier on the input.
  This entry says that the definition exists, not that it holds for any given $c$.
  \own{}
  \leanfile{Sagitta/Map/Borde/Goodstein.lean}
\end{definition}

\begin{theorem}[Satisfiability on a polynomially decidable set of lengths is in $\cNP$]
  \label{thm:structure-sat-en-np}
  \lean{Sagitta.Map.Borde.Goodstein.satEn_NP}
  \leanok
  \uses{def:structure-sat-en, def:structure-longitudes-polinomiales, def:NP}
  Let $c : \mathbb{N} \to \mathrm{Bool}$ and assume $\mathrm{LongitudesPolinomiales}\,c$. Then
  $\mathrm{satEn}\,c \in \cNP$ for the encoding of CNF formulas.
  \own{}
  \leanfile{Sagitta/Map/Borde/Goodstein.lean}
\end{theorem}

\begin{proof}
  \leanok
  Take the verifier of $\SAT$ from cslib-frontier and put the length test in front of it
  (as a guard); both are computable in polynomial time, and so is their conjunction.
\end{proof}

\section{Pumping: the lemma and the move against $\SAT$}
\label{sec:structure-pumping}

A pumping lemma has two halves: what every machine of the model satisfies, and a word that breaks it.
The statements of this section state the break as a definition (the \emph{move}) and prove that, for
search problems solved by the Turing machines of cslib, the move is the negation of the lemma; the
move against $\SAT$ itself is not constructed here. The second group is the finite-memory case, where the move
can be exhibited, and the third is a count in an abstract model of rivals.

\begin{definition}[Depth move for a search problem]
  \label{def:structure-jugada-profundidad}
  \lean{Sagitta.Map.Borde.PBombeo.JugadaProfundidad, Sagitta.Map.Tiempo.LevinMaquinas.KtMTLog, Sagitta.Map.Tiempo.LevinMaquinas.ktMT, Sagitta.Map.Tiempo.LevinMaquinas.ResolubleMT}
  \leanok
  Let $\enc : \alpha \hookrightarrow \mathrm{List}\,\mathrm{Bool}$ be an encoding, $V : \alpha
  \times \mathrm{List}\,\mathrm{Bool} \to \mathrm{Bool}$ a verifier and $A \subseteq \alpha$. An
  instance $a$ is solvable when some machine of cslib, run for some number of steps, produces a
  witness accepted by $V$; $\mathrm{kt}(a)$ is the least value of $i + \lceil \log_2 t \rceil$
  over the machines $i$ and step counts $t$ with which $a$ is solved ($0$ if there is none). The
  property $\mathrm{KtMTLog}$ of $A$ says that for some $c$ every $a \in A$ is solvable with
  $\mathrm{kt}(a) \le c \lfloor \log_2 (|\enc(a)| + 2) \rfloor$. The \emph{depth move} says that for
  every $c$ there is $a \in A$ such that, if $a$ is solvable, $c \lfloor \log_2 (|\enc(a)| + 2)
  \rfloor < \mathrm{kt}(a)$.
  \own{}
  \leanfile{Sagitta/Map/Borde/PBombeo.lean}
\end{definition}

\begin{theorem}[The depth move is the negation of bounded depth]
  \label{thm:structure-jugada-iff}
  \lean{Sagitta.Map.Borde.PBombeo.jugada_iff_no_bombea}
  \leanok
  \uses{def:structure-jugada-profundidad}
  Let $\enc$, $V$ and $A$ be as in the definition. The depth move holds for $(\enc, V, A)$ if and
  only if $\mathrm{KtMTLog}$ does not hold for $A$.
  \own{}
  \leanfile{Sagitta/Map/Borde/PBombeo.lean}
\end{theorem}

\begin{proof}
  \leanok
  Unfold both sides and push the negation through the quantifiers.
\end{proof}

\begin{definition}[Finite-memory move]
  \label{def:structure-jugada-finita}
  \lean{Sagitta.Map.Borde.BombeoFinito.JugadaFinita, Sagitta.Map.Borde.BombeoFinito.Distinguibles, Sagitta.Map.Borde.BombeoFinito.anbn}
  \leanok
  Two words $u, v$ over an alphabet $\Gamma$ are distinguished by a language $L$ when the left
  quotients of $L$ by $u$ and by $v$ differ. The \emph{finite-memory move} for $L$ says that for
  every $k$ there are $k + 1$ words, indexed by $\mathrm{Fin}(k+1)$, any two of which (with different
  indices) are distinguished by $L$. For letters $a, b$, the language $\mathrm{anbn}$ is the set of
  words $a^n b^n$.
  \own{}
  \leanfile{Sagitta/Map/Borde/BombeoFinito.lean}
\end{definition}

\begin{theorem}[Pumping for finite automata, from the pigeonhole principle]
  \label{thm:structure-bombeo-finito}
  \lean{Sagitta.Map.Borde.BombeoFinito.bombeo_finito}
  \leanok
  Let $M$ be a deterministic finite automaton over $\Gamma$ with a finite set $\sigma$ of states,
  $q \in \sigma$ a state and $w$ a word with $|\sigma| \le |w|$. Then there are words $x, y, z$ with
  $w = xyz$, $y \neq \varepsilon$, $|x| + |y| \le |\sigma|$, and such that for every $m$ the
  automaton started at $q$ ends in the same state on $x y^m z$ as on $w$.
  \own{}
  \leanfile{Sagitta/Map/Borde/BombeoFinito.lean}
\end{theorem}

\begin{proof}
  \leanok
  Among the $|\sigma| + 1$ prefixes of $w$ of length at most $|\sigma|$, two lead to the same state
  (pigeonhole principle); the factor between them is $y$, and induction on $m$ shows that
  repeating it does not change the final state.
\end{proof}

\begin{theorem}[The move for $\{a^n b^n\}$ is exhibited]
  \label{thm:structure-anbn-jugada}
  \lean{Sagitta.Map.Borde.BombeoFinito.anbn_jugada}
  \leanok
  \uses{def:structure-jugada-finita}
  For letters $a \neq b$ of $\Gamma$, the finite-memory move holds for $\mathrm{anbn}(a, b)$.
  \own{}
  \leanfile{Sagitta/Map/Borde/BombeoFinito.lean}
\end{theorem}

\begin{proof}
  \leanok
  The words $a^0, a^1, \dots, a^k$ are distinguished by the suffix $b^i$, which $\mathrm{anbn}$
  accepts after $a^i$ and after no other power of $a$.
\end{proof}

\begin{theorem}[The depth beaten by a diagonal instance is at most $\log_2$ of its formula size, in an abstract count]
  \label{thm:structure-profundidad-normalizada}
  \lean{Sagitta.Map.Borde.JugadaDiagonal.profundidad_normalizada}
  \leanok
  Let $\ell, t, N$ be natural numbers with $0 < t$ and $2^\ell \cdot t \le N$. Then
  $\ell + \lfloor \log_2 t \rfloor \le \lfloor \log_2 N \rfloor$. There is no machine in the
  statement: $2^\ell$ stands for the number of rivals of level $\ell$, $t$ for the steps of each, and
  $N$ for the size of a formula that writes all the runs.
  \own{}
  \leanfile{Sagitta/Map/Borde/JugadaDiagonal.lean}
\end{theorem}

\begin{proof}
  \leanok
  $2^{\ell + \lfloor \log_2 t \rfloor} = 2^\ell \cdot 2^{\lfloor \log_2 t \rfloor} \le 2^\ell \cdot t
  \le N$, and then the characterization of the floor of the logarithm.
\end{proof}

\section{Lower bounds on backtracking and on cutting planes}
\label{sec:structure-bounds}

The statements of this section are unconditional lower bounds in two restricted models. The first
is a backtracking search over truth assignments of clauses, in an abstract cost model that counts
the nodes visited; the second is the proof system of cutting planes, where a refutation is a list of
integer linear inequalities.

\begin{definition}[Backtracking on clauses]
  \label{def:structure-backtracking}
  \lean{Sagitta.Map.Tiempo.Vivo, Sagitta.Map.Tiempo.nodos, Sagitta.Map.Tiempo.buscar, Sagitta.Map.Tiempo.solCNF, Sagitta.Map.Tiempo.detCNF}
  \leanok
  A clause is a list of literals, a literal being a variable index with a sign. A node is a list $p$
  of Booleans (a partial assignment of the first $|p|$ variables), and $k$ variables remain to be
  fixed. The node is \emph{alive} for a test $\mathrm{sol}$ on lists of Booleans when some list
  $q$ of length $k$ has $\mathrm{sol}(p \mathbin{+\!\!+} q) = \mathrm{true}$. For a list $F$ of
  clauses, $\mathrm{solCNF}(F)$ accepts an assignment that makes every clause true, and the
  detector $\mathrm{detCNF}(F)$ says that some clause of $F$ is already false in $p$ (all its
  variables are fixed and none with its sign). The search $\mathrm{buscar}$ returns at a node whose
  detector says dead; otherwise it tries $p\,\mathtt{false}$ and, if that branch finds nothing,
  $p\,\mathtt{true}$; $\mathrm{nodos}$ is the number of nodes it visits.
  \own{}
  \leanfile{Sagitta/Map/ArbolElecciones.lean}
\end{definition}

\begin{lemma}[An unsatisfiable formula has a dead root]
  \label{lem:structure-raiz-muerta}
  \lean{Sagitta.Map.Tiempo.raiz_muerta}
  \leanok
  \uses{def:structure-backtracking}
  Let $F$ be a list of clauses, read as a CNF formula over $\mathbb{N}$, that is not satisfiable,
  and let $N$ be a natural number. Then the root (the empty partial assignment, with $N$ variables
  to fix) is not alive for $\mathrm{solCNF}(F)$.
  \source{\cite{BeyersdorffGalesiLauria13}.}
  \fidelity{Not checked against the printed source.}
  \leanfile{Sagitta/Map/RaizMuerta.lean}
\end{lemma}

\begin{proof}
  \leanok
  A completion of the root accepted by $\mathrm{solCNF}(F)$ is read as an assignment that satisfies
  every clause, which contradicts the hypothesis.
\end{proof}

\begin{theorem}[Backtracking with the falsified-clause detector visits many nodes on the pigeonhole formula]
  \label{thm:structure-palomar-nodos}
  \lean{Sagitta.Map.Tiempo.palomar_nodos}
  \leanok
  \uses{def:structure-backtracking, lem:structure-raiz-muerta}
  Let $m, n, h$ be natural numbers with $n < m$ and $0 < n$, let $f : \mathrm{Fin}\,m \times
  \mathrm{Fin}\,n \to \mathbb{N}$ be injective, and let $N$ be such that $f(v) < N$ for every $v$.
  Let $F$ be the pigeonhole formula $\mathrm{PHP}^m_n$ with each variable $v$ renamed to $f(v)$.
  Then $(n + 1 - h)^h < 2 \cdot \mathrm{nodos}$, where $\mathrm{nodos}$ is the number of nodes that
  backtracking with the detector $\mathrm{detCNF}(F)$ visits from the root with $N$ variables to
  fix (the subtraction $n + 1 - h$ is that of natural numbers).
  \source{\cite{BeyersdorffGalesiLauria13}.}
  \fidelity{Not checked against the printed source.}
  \leanfile{Sagitta/Map/RaizMuerta.lean}
\end{theorem}

\begin{proof}
  \leanok
  The run of the search is a decision tree that finds a falsified clause (it is a leaf at every
  dead node, and the root is dead by Lemma~\ref{lem:structure-raiz-muerta}); renaming the
  variables back with the inverse of $f$ gives a tree for $\mathrm{PHP}^m_n$ with the same leaves.
  The bound of cslib-frontier on the leaves of such trees applies, and the number of leaves is at
  most the number of nodes visited.
\end{proof}

\begin{definition}[Cutting-plane refutation with bounded coefficients]
  \label{def:structure-refuta-acotada}
  \lean{Sagitta.Map.Tiempo.PlanosDeCorteClique.RefutaAcotada}
  \leanok
  For a set $\Gamma$ of integer linear inequalities and natural numbers $s, M$, $\Gamma$ has a
  cutting-plane refutation of at most $s$ lines with coefficients of absolute value at most $M$
  when some list $ls$ of inequalities is a cutting-plane proof from $\Gamma$, contains the
  contradiction, has length at most $s$, and every coefficient of every inequality of $ls$ has
  absolute value at most $M$.
  \own{}
  \leanfile{Sagitta/Map/PlanosDeCorteClique.lean}
\end{definition}

\begin{theorem}[No polynomial cutting-plane refutations with polynomial coefficients of the clique-coloring formulas]
  \label{thm:structure-sin-refutacion-polinomial}
  \lean{Sagitta.Map.Tiempo.PlanosDeCorteClique.sin_refutacion_polinomial}
  \leanok
  \uses{def:structure-refuta-acotada}
  There are no functions $p, q : \mathbb{N} \to \mathbb{N}$ bounded by polynomials such that, for
  infinitely many $j$, the inequalities of the clique-coloring formula with $n = 2^{4j}$ vertices,
  clique size $k = 2^{2j} + 1$ and $r = 2^{2j}$ colors have a cutting-plane refutation of at most
  $p(2^{4j})$ lines with coefficients of absolute value at most $q(2^{4j})$.
  \source{\cite{Pudlak97,Krajicek97,AlonBoppana87}.}
  \fidelity{Not checked against the printed source.}
  \leanfile{Sagitta/Map/PlanosDeCorteClique.lean}
\end{theorem}

\begin{proof}
  \leanok
  A refutation with bounded coefficients gives, by the interpolation theorem of cslib-frontier, a
  monotone circuit that separates the graphs with a $k$-clique from the $r$-colorable ones, of size
  at most $(s \cdot 2^{2t} + 1)\, 2^{2t+1}$ with $t$ determined by the coefficient bound. The monotone
  lower bound of cslib-frontier makes this size exceed $2^{\lfloor \lfloor \sqrt{\lfloor \sqrt n
  \rfloor} \rfloor / 2 \rfloor}$, and for a large enough $j$ the polynomial bounds contradict this.
\end{proof}

\section{Grey boxes: resolution with bounded memory, a ladder of proof systems, and the limit}
\label{sec:structure-grey}

A \emph{grey box} sees the formula but works only with fixed rules. The statements of this section
live in three models: resolution on configurations of clauses with a bound on the number of clauses
held in memory; an abstract ladder of proof systems, each given by the sizes of the proofs of an
indexed family of instances; and proof systems in the sense of Cook and Reckhow for the
unsatisfiable CNF formulas.

\begin{definition}[Configurational resolution in bounded space]
  \label{def:structure-caja-gris}
  \lean{Sagitta.Map.Tiempo.CajaGris.Jugada, Sagitta.Map.Tiempo.CajaGris.Avanza, Sagitta.Map.Tiempo.CajaGris.Recorre, Sagitta.Map.Tiempo.CajaGris.largoDe, Sagitta.Map.Tiempo.CajaGris.RefutaEn, Sagitta.Map.Tiempo.CajaGris.numVars}
  \leanok
  A configuration is a finite set of clauses. A move loads an axiom of the formula $F$ into the
  configuration, adds the resolvent of two clauses of the configuration on a variable (that
  occurs positively in the first and negatively in the second), or deletes a clause. A run in space
  $s$ is a list of moves from the empty configuration in which every configuration reached has at
  most $s$ clauses; its length counts loads and resolutions, not deletions. $F$ has a refutation in
  space $s$ and length at most $L$ when some run in space $s$ of length at most $L$ ends in a
  configuration that contains the empty clause. $\mathrm{numVars}(c)$ is the number of distinct
  variables of the clause $c$.
  \own{}
  \leanfile{Sagitta/Map/CajaGris.lean}
\end{definition}

\begin{theorem}[Space is at least the minimal width of the axioms plus one]
  \label{thm:structure-espacio-ancho}
  \lean{Sagitta.Map.Tiempo.CajaGris.espacio_ge_ancho}
  \leanok
  \uses{def:structure-caja-gris}
  Let $F$ be a list of clauses and $s, L, w$ natural numbers. If $F$ has a refutation in space $s$
  and length at most $L$, and every axiom of $F$ has at least $w$ distinct variables, then $w + 1 \le
  s$.
  \own{}
  \leanfile{Sagitta/Map/CajaGris.lean}
\end{theorem}

\begin{proof}
  \leanok
  Follow the run to the first load of an axiom $A$ that turns a satisfiable configuration $N$ into an
  unsatisfiable one (resolving and deleting do not remove solutions). Flipping each variable of $A$ in
  a model of $N$ satisfies $A$ and falsifies a different clause of $N$ for each variable, so $N$ has
  at least $w$ clauses and the configuration after the load has at least $w + 1$.
\end{proof}

\begin{definition}[A ladder of proof systems over a family of instances]
  \label{def:structure-peldano}
  \lean{Sagitta.Map.Borde.EscaleraGris.Peldano, Sagitta.Map.Borde.EscaleraGris.Conserva, Sagitta.Map.Borde.EscaleraGris.Corto, Sagitta.Map.Borde.EscaleraGris.Simula, Sagitta.Map.Borde.EscaleraGris.arbol, Sagitta.Map.Borde.EscaleraGris.resolucion, Sagitta.Map.Borde.EscaleraGris.planos}
  \leanok
  A \emph{rung} over a type $\iota$ of instances is a predicate $R(i, s)$ (``the instance $i$ has a
  proof of size at most $s$''). It \emph{conserves} the floor $c : \iota \to \mathbb{N}$ when $R(i, s)$
  implies $c(i) \le s$; it is \emph{short} with ceiling $a$ when $R(i, a(i))$ holds for every $i$;
  $Q$ \emph{simulates} $P$ with overhead $f$ when $P(i, s)$ implies $Q(i, f(i, s))$. Over the
  pigeonhole formulas $\mathrm{PHP}^{n+1}_n$ the rungs $\mathrm{arbol}$, $\mathrm{resolucion}$ and
  $\mathrm{planos}$ say that the formula has a tree-like resolution refutation of at most $s$
  nodes, a resolution refutation of at most $s$ lines, and a cutting-plane refutation of at most $s$
  lines.
  \own{}
  \leanfile{Sagitta/Map/Borde/EscaleraGris.lean}
\end{definition}

\begin{theorem}[Resolution cannot simulate cutting planes on the pigeonhole formulas with small overhead]
  \label{thm:structure-separa-resolucion-planos}
  \lean{Sagitta.Map.Borde.EscaleraGris.separa_resolucion_planos}
  \leanok
  \uses{def:structure-peldano}
  Let $f : \mathbb{N} \to \mathbb{N} \to \mathbb{N}$ and assume that $\mathrm{resolucion}$ simulates
  $\mathrm{planos}$ with overhead $f$ over the family $\mathrm{PHP}^{n+1}_n$. Then for every $n$,
  $2^{\lfloor n/24 \rfloor} \le f(n, (n+1)(n+2)^2)$.
  \own{}
  \leanfile{Sagitta/Map/Borde/EscaleraGris.lean}
\end{theorem}

\begin{proof}
  \leanok
  Cutting planes refute $\mathrm{PHP}^{n+1}_n$ in $(n+1)(n+2)^2$ lines (it is short), resolution
  conserves the floor $2^{\lfloor n/24 \rfloor}$ (the resolution lower bound of cslib-frontier), and an
  abstract lemma says that a simulation of a short rung by a rung that conserves a floor pays that
  floor.
\end{proof}

\begin{definition}[Hypothesis: conservation holds in every grey box]
  \label{def:structure-ley-caja-gris}
  \lean{Sagitta.Map.Borde.LimiteBlanco.LeyEnTodaCajaGris, Sagitta.Map.Borde.LimiteBlanco.ConservaCR, Sagitta.Map.Borde.LimiteBlanco.NPigualCoNP}
  \leanok
  \uses{def:P, def:NP}
  $\mathrm{NPigualCoNP}$ says that $\cNP = \cconP$ for every type and every encoding into binary
  words. For a proof system $P$ for the unsatisfiable CNF formulas, in the sense of Cook and
  Reckhow (cslib-frontier), $\mathrm{ConservaCR}(P)$ says that $P$ is not polynomially bounded, that
  is, that no polynomial bounds the length of a shortest proof of each unsatisfiable formula.
  $\mathrm{LeyEnTodaCajaGris}$ says that every proof system for the unsatisfiable CNF formulas
  satisfies $\mathrm{ConservaCR}$. This entry says that the definitions exist, not that the law holds;
  it is open.
  \own{}
  \leanfile{Sagitta/Map/Borde/LimiteBlanco.lean}
\end{definition}

\begin{theorem}[The law in every grey box is $\cNP \neq \cconP$]
  \label{thm:structure-ley-iff}
  \lean{Sagitta.Map.Borde.LimiteBlanco.ley_iff_np_ne_conp}
  \leanok
  \uses{def:structure-ley-caja-gris}
  $\mathrm{LeyEnTodaCajaGris}$ holds if and only if $\mathrm{NPigualCoNP}$ does not hold.
  \own{}
  \leanfile{Sagitta/Map/Borde/LimiteBlanco.lean}
\end{theorem}

\begin{proof}
  \leanok
  Unfold the definitions and rewrite with the theorem of cslib-frontier that $\cNP = \cconP$ if and
  only if some proof system is polynomially bounded; the negation is moved outside the existential.
\end{proof}

\begin{theorem}[The law in every grey box would imply $\cP \neq \cNP$]
  \label{thm:structure-ley-implica}
  \lean{Sagitta.Map.Borde.LimiteBlanco.ley_implica_P_ne_NP}
  \leanok
  \uses{def:structure-ley-caja-gris, thm:structure-ley-iff, def:P, def:NP}
  Assume $\mathrm{LeyEnTodaCajaGris}$. Then it is not the case that $\cP = \cNP$ for every type and
  every encoding into binary words.
  \own{}
  \leanfile{Sagitta/Map/Borde/LimiteBlanco.lean}
\end{theorem}

\begin{proof}
  \leanok
  If $\cP = \cNP$ for every encoding, then $\cNP$ is closed under complement, so $\cNP = \cconP$,
  which contradicts Theorem~\ref{thm:structure-ley-iff}.
\end{proof}

\section{The approximation rate and the Peano ceiling}
\label{sec:structure-rate}

The statements of this section are about an enumeration $M_0, M_1, \dots$ of problems, that is,
functions from binary words to Booleans (the problems of the map), and a problem $L$. No machine
appears: the enumeration is a parameter.

\begin{definition}[Approximation rate]
  \label{def:structure-tasa}
  \lean{Sagitta.Map.PrimerFallo.FallaEn, Sagitta.Map.PrimerFallo.Falla, Sagitta.Map.PrimerFallo.primerFallo, Sagitta.Map.PrimerFallo.tasa}
  \leanok
  Machine $j$ \emph{fails at length $n$} when some word $x$ of length $n$ has $M_j(x) \neq L(x)$, and
  it fails when it fails at some length. The first failure of $j$ is the least such length ($0$ if it
  never fails). The approximation rate $R(i)$ of $L$ by the enumeration is the partial function whose
  value is the maximum of the first failures of the machines $0, \dots, i - 1$, defined only when all
  of them fail.
  \source{\cite{BenDavidHalevi91}.}
  \fidelity{Not checked against the printed source.}
  \leanfile{Sagitta/Map/PrimerFallo.lean}
\end{definition}

\begin{lemma}[The rate is total if and only if $L$ is not in the enumeration]
  \label{lem:structure-tasa-total}
  \lean{Sagitta.Map.PrimerFallo.tasa_total_iff}
  \leanok
  \uses{def:structure-tasa}
  Let $M : \mathbb{N} \to \mathrm{Problema}$ and $L$ be a problem. The rate $R(i)$ is defined for every $i$
  if and only if $L$ is not in the range of $M$.
  \source{\cite{BenDavidHalevi91}.}
  \fidelity{Not checked against the printed source.}
  \leanfile{Sagitta/Map/PrimerFallo.lean}
\end{lemma}

\begin{proof}
  \leanok
  Failing at some input is the same as being different from $L$; if $L = M_j$ then $R(j + 1)$ is
  undefined, and if $L$ is none of them every machine fails.
\end{proof}

\begin{definition}[Hypothesis: $\cP \subseteq \cNP$ in the class of the map]
  \label{def:structure-psubnp}
  \lean{Sagitta.Map.PSubNP, Sagitta.Map.PDistintoNP}
  \leanok
  \uses{def:structure-tasa}
  % uses-lean: Sagitta.Map.P, Sagitta.Map.NP
  $\mathrm{PSubNP}$ says that the class $\cP$ of the map is contained in its class $\cNP$;
  $\mathrm{PDistintoNP}$ says that the two classes are different. This entry says that the
  definitions exist, not that $\cP \subseteq \cNP$ is assumed to hold.
  \own{}
  \leanfile{Sagitta/Map/Clases.lean}
\end{definition}

\begin{theorem}[$\cP \neq \cNP$ in terms of the approximation rate]
  \label{thm:structure-pnp-tasa}
  \lean{Sagitta.Map.PrimerFallo.pDistintoNP_iff_tasa_total}
  \leanok
  \uses{def:structure-tasa, def:structure-psubnp, lem:structure-tasa-total}
  % uses-lean: Sagitta.Map.P, Sagitta.Map.NP
  Let $M : \mathbb{N} \to \mathrm{Problema}$ be an enumeration whose range is exactly the class $\cP$ of
  the map, and assume $\mathrm{PSubNP}$. Then $\cP \neq \cNP$ if and only if there is $L$ in the class
  $\cNP$ of the map such that $R(i)$ is defined for every $i$.
  \source{\cite{BenDavidHalevi91}.}
  \fidelity{Not checked against the printed source.}
  \leanfile{Sagitta/Map/PrimerFallo.lean}
\end{theorem}

\begin{proof}
  \leanok
  Totality of the rate is $L \notin \cP$ by the lemma for $\cP$; with $\cP \subseteq \cNP$,
  the classes differ exactly when some member of $\cNP$ is outside $\cP$.
\end{proof}

\begin{definition}[Statement (not proved): the Peano ceiling on the approximation rate]
  \label{def:structure-teorema4}
  \lean{Sagitta.Map.Techo.Teorema4, Sagitta.Map.Techo.TasaBajoElTecho, Sagitta.Map.Techo.Domina}
  \leanok
  \uses{def:structure-tasa}
  A function $f$ dominates $g$ when $g(n) < f(n)$ for all large $n$. The rate $R$ of $L$ by $M$ is
  \emph{under the ceiling} when there is a notation $o$ in Cantor normal form (an ordinal below
  $\varepsilon_0$) and an $N$ such that for every $i \ge N$ the rate $R(i)$ is defined and lies below
  $F_o(i)$, with $F_o$ the fast-growing hierarchy of Mathlib. For a proposition $\mathrm{Demuestra}$,
  $\mathrm{Teorema4}$ is the statement ``$\mathrm{Demuestra}$ if and only if the rate is under the
  ceiling''. There is no theory in Lean: $\mathrm{Demuestra}$ is a parameter standing for ``the theory
  proves that $L$ is not in the enumeration'', and this entry says that the statement is defined, not
  that it holds.
  \source{\cite{BenDavidHalevi91}.}
  \fidelity{Not checked against the printed source.}
  \leanfile{Sagitta/Map/Techo.lean}
\end{definition}

\section{Dichotomy against intermediate problems: the three switches}
\label{sec:structure-switches}

The statements of this section are about classes of languages over binary words, with the classes
$\cP$ and $\cNP$, the polynomial-time many-one reductions and the $\cNP$-completeness of
cslib-frontier. A class is a set of languages. The syntax of the fragments of existential
second-order logic is not formalized: each fragment is a parameter, a class of languages, and what is
taken from the literature about it is gathered in a named hypothesis.

\begin{definition}[Polynomial equivalence, dichotomy, intermediate problems, capture of $\cNP$]
  \label{def:structure-dicotomia-defs}
  \lean{Sagitta.Map.Borde.TresInterruptores.Lenguaje, Sagitta.Map.Borde.TresInterruptores.Clase, Sagitta.Map.Borde.TresInterruptores.Reduce, Sagitta.Map.Borde.TresInterruptores.Equivalentes, Sagitta.Map.Borde.TresInterruptores.TieneDicotomia, Sagitta.Map.Borde.TresInterruptores.TieneIntermedios, Sagitta.Map.Borde.TresInterruptores.CapturaNP}
  \leanok
  \uses{def:reduction, def:P, def:NP}
  A language is a set of binary words and a class is a set of languages. $\mathrm{Reduce}(L, L')$ is
  the polynomial-time many-one reduction $L \red L'$ for the binary encoding of words;
  $L$ and $L'$ are \emph{equivalent} when both $L \red L'$ and $L' \red L$. A class $F$ \emph{has the
  dichotomy} when every $L \in F$ is in $\cP$ or $\cNP$-complete. It \emph{has intermediate problems}
  when some $L \in F$ is in $\cNP$, not in $\cP$ and not $\cNP$-complete. It \emph{captures} $\cNP$
  when every $L \in \cNP$ is equivalent to some $L' \in F$.
  \own{}
  \leanfile{Sagitta/Map/Borde/TresInterruptores.lean}
\end{definition}

\begin{theorem}[Being intermediate is preserved by polynomial equivalence]
  \label{thm:structure-intermedio-equiv}
  \lean{Sagitta.Map.Borde.TresInterruptores.Equivalentes.intermedio_iff}
  \leanok
  \uses{def:structure-dicotomia-defs, def:structure-intermedio}
  If $L$ and $L'$ are equivalent, then $L$ is intermediate if and only if $L'$ is.
  \own{}
  \leanfile{Sagitta/Map/Borde/TresInterruptores.lean}
\end{theorem}

\begin{proof}
  \leanok
  Membership in $\cNP$, membership in $\cP$ and $\cNP$-completeness are each preserved by
  equivalence: the first two go down along a reduction, and completeness goes up along a reduction to
  a language in $\cNP$.
\end{proof}

\begin{theorem}[The dichotomy is transported by polynomial equivalence]
  \label{thm:structure-dicotomia-equiv}
  \lean{Sagitta.Map.Borde.TresInterruptores.dicotomia_de_equivalentes}
  \leanok
  \uses{def:structure-dicotomia-defs}
  Let $F, G$ be classes and assume that $G$ has the dichotomy and that every $L \in F$ is
  equivalent to some $C \in G$. Then $F$ has the dichotomy.
  \own{}
  \leanfile{Sagitta/Map/Borde/TresInterruptores.lean}
\end{theorem}

\begin{proof}
  \leanok
  If $C$ is in $\cP$ then so is $L$, by $L \red C$; if $C$ is $\cNP$-complete then $L \in \cNP$ by
  $L \red C$ and $L$ is hard by $C \red L$.
\end{proof}

\begin{theorem}[In a class that captures $\cNP$, intermediate problems exist if and only if $\cP \neq \cNP$]
  \label{thm:structure-captura-iff}
  \lean{Sagitta.Map.Borde.TresInterruptores.captura_iff, Sagitta.Map.Borde.TresInterruptores.intermedios_de_captura, Sagitta.Map.Borde.TresInterruptores.P_ne_NP_de_intermedios}
  \leanok
  \uses{def:structure-dicotomia-defs, thm:structure-ladner-fractal, thm:structure-intermedio-equiv}
  Let $F$ be a class that captures $\cNP$. Then $F$ has intermediate problems if and only if $\cP \neq
  \cNP$ over binary words. The direction from $\cP \neq \cNP$ is \texttt{intermedios\_de\_captura}; the
  converse holds for every class, whether or not it captures $\cNP$ (\texttt{P\_ne\_NP\_de\_intermedios}).
  \own{}
  \leanfile{Sagitta/Map/Borde/TresInterruptores.lean}
\end{theorem}

\begin{proof}
  \leanok
  An intermediate problem is in $\cNP$ and not in $\cP$. Conversely, Theorem~\ref{thm:structure-ladner-fractal}
  gives an intermediate language, which is equivalent to some language of $F$, and being intermediate
  is preserved by equivalence.
\end{proof}

\begin{theorem}[If $\cP = \cNP$, a language of $\cNP$ is complete exactly when it is non-trivial]
  \label{thm:structure-completo-no-trivial}
  \lean{Sagitta.Map.Borde.TresInterruptores.completo_iff_no_trivial}
  \leanok
  \uses{thm:structure-completo-de-p-eq-np}
  Assume that $\cP = \cNP$ for every type and every encoding into binary words. For a language $L
  \in \cNP$ over binary words, $L$ is $\cNP$-complete if and only if $L \neq \emptyset$ and $L \neq
  \Sigma^*$.
  \own{}
  \leanfile{Sagitta/Map/Borde/TresInterruptores.lean}
\end{theorem}

\begin{proof}
  \leanok
  Neither the empty language nor the set of all words is $\cNP$-complete, since the other one is in
  $\cNP$ and does not reduce to it. A language with a member and a non-member is complete by
  Theorem~\ref{thm:structure-completo-de-p-eq-np}.
\end{proof}

\begin{definition}[The three switches]
  \label{def:structure-interruptores}
  \lean{Sagitta.Map.Borde.TresInterruptores.Interruptores, Sagitta.Map.Borde.TresInterruptores.encendido, Sagitta.Map.Borde.TresInterruptores.sinM, Sagitta.Map.Borde.TresInterruptores.sinMo, Sagitta.Map.Borde.TresInterruptores.sinD, Sagitta.Map.Borde.TresInterruptores.Impone}
  \leanok
  A setting of the switches is a triple of Booleans (monotone, monadic, without inequality), where
  true means that the restriction is imposed. $\mathrm{encendido}$ is the setting with all three
  imposed; $\mathrm{sinM}$, $\mathrm{sinMo}$ and $\mathrm{sinD}$ are the settings in which only the
  first, only the second, only the third is switched off. $\mathrm{Impone}(s, t)$ says that every
  restriction imposed by $t$ is imposed by $s$.
  \own{}
  \leanfile{Sagitta/Map/Borde/TresInterruptores.lean}
\end{definition}

\begin{definition}[Hypothesis: what is taken from Feder and Vardi]
  \label{def:structure-feder-vardi}
  \lean{Sagitta.Map.Borde.TresInterruptores.FederVardi}
  \leanok
  \uses{def:structure-interruptores, def:structure-dicotomia-defs}
  Let $\mathrm{frag}$ assign a class of languages to each setting of the switches and let $\mathrm{csp}$
  be a class of languages. $\mathrm{FederVardi}(\mathrm{frag}, \mathrm{csp})$ is the conjunction of
  the following statements.
  \begin{itemize}
    \item If $\mathrm{Impone}(s, t)$ then $\mathrm{frag}(s) \subseteq \mathrm{frag}(t)$.
    \item $\mathrm{frag}(s) \subseteq \cNP$ for every $s$.
    \item Every $L \in \mathrm{frag}(\mathrm{encendido})$ is equivalent to some $C \in \mathrm{csp}$.
    \item $\mathrm{frag}(\mathrm{sinM})$, $\mathrm{frag}(\mathrm{sinMo})$ and $\mathrm{frag}(\mathrm{sinD})$
      each capture $\cNP$.
  \end{itemize}
  The hypothesis is a structure of type \texttt{Prop}; this entry says that it is defined, not that it
  holds.
  \own{}
  \leanfile{Sagitta/Map/Borde/TresInterruptores.lean}
\end{definition}

\begin{theorem}[With one switch off, the fragment captures $\cNP$]
  \label{thm:structure-captura-apagado}
  \lean{Sagitta.Map.Borde.TresInterruptores.captura_de_apagado}
  \leanok
  \uses{def:structure-feder-vardi}
  Assume $\mathrm{FederVardi}(\mathrm{frag}, \mathrm{csp})$ and let $s$ be a setting with $s \neq
  \mathrm{encendido}$. Then $\mathrm{frag}(s)$ captures $\cNP$.
  \own{}
  \leanfile{Sagitta/Map/Borde/TresInterruptores.lean}
\end{theorem}

\begin{proof}
  \leanok
  Case analysis on the three switches: the setting imposes at least the restrictions of one of the
  three corners, so its fragment is contained in the fragment of that corner, and capturing $\cNP$ is
  inherited by larger classes.
\end{proof}

\begin{theorem}[All three switches on: the dichotomy, given the dichotomy of the finite-template class]
  \label{thm:structure-encendido-dicotomia}
  \lean{Sagitta.Map.Borde.TresInterruptores.encendido_dicotomia}
  \leanok
  \uses{def:structure-feder-vardi, thm:structure-dicotomia-equiv}
  Assume $\mathrm{FederVardi}(\mathrm{frag}, \mathrm{csp})$ and that the class $\mathrm{csp}$ has the
  dichotomy. Then $\mathrm{frag}(\mathrm{encendido})$ has the dichotomy.
  \own{}
  \leanfile{Sagitta/Map/Borde/TresInterruptores.lean}
\end{theorem}

\begin{proof}
  \leanok
  Every language of the fragment is equivalent to a member of $\mathrm{csp}$
  (Theorem~\ref{thm:structure-dicotomia-equiv}).
\end{proof}

\begin{theorem}[With a switch off, intermediate problems exist if and only if $\cP \neq \cNP$]
  \label{thm:structure-apagado-iff}
  \lean{Sagitta.Map.Borde.TresInterruptores.apagado_iff}
  \leanok
  \uses{def:structure-feder-vardi, thm:structure-captura-apagado, thm:structure-captura-iff}
  Assume $\mathrm{FederVardi}(\mathrm{frag}, \mathrm{csp})$ and $s \neq \mathrm{encendido}$. Then
  $\mathrm{frag}(s)$ has intermediate problems if and only if $\cP \neq \cNP$ over binary words.
  \own{}
  \leanfile{Sagitta/Map/Borde/TresInterruptores.lean}
\end{theorem}

\begin{proof}
  \leanok
  Theorem~\ref{thm:structure-captura-iff} applied to the class $\mathrm{frag}(s)$, which captures
  $\cNP$ by Theorem~\ref{thm:structure-captura-apagado}.
\end{proof}

\begin{theorem}[The switch decides which side of the table a fragment is on]
  \label{thm:structure-interruptor-decide}
  \lean{Sagitta.Map.Borde.TresInterruptores.interruptor_decide, Sagitta.Map.Borde.TresInterruptores.dicotomia_iff_encendido}
  \leanok
  \uses{def:structure-feder-vardi, thm:structure-encendido-dicotomia, thm:structure-apagado-iff}
  Assume $\mathrm{FederVardi}(\mathrm{frag}, \mathrm{csp})$ and that $\mathrm{csp}$ has the dichotomy.
  For every setting $s$, either $s = \mathrm{encendido}$ and $\mathrm{frag}(s)$ has the dichotomy, or $s
  \neq \mathrm{encendido}$ and $\mathrm{frag}(s)$ has intermediate problems if and only if $\cP \neq
  \cNP$ (\texttt{interruptor\_decide}). If moreover $\cP \neq \cNP$ over binary words, then
  $\mathrm{frag}(s)$ has the dichotomy if and only if $s = \mathrm{encendido}$
  (\texttt{dicotomia\_iff\_encendido}).
  \own{}
  \leanfile{Sagitta/Map/Borde/TresInterruptores.lean}
\end{theorem}

\begin{proof}
  \leanok
  Split on $s = \mathrm{encendido}$ and use the two previous theorems; for the second statement a
  fragment with a switch off that had the dichotomy could not have an intermediate problem, but
  $\cP \neq \cNP$ gives one.
\end{proof}

\begin{theorem}[If $\cP = \cNP$, every fragment has the dichotomy and none has intermediate problems]
  \label{thm:structure-todos-dicotomia}
  \lean{Sagitta.Map.Borde.TresInterruptores.todos_dicotomia_de_P_eq_NP}
  \leanok
  \uses{def:structure-feder-vardi}
  Assume $\mathrm{FederVardi}(\mathrm{frag}, \mathrm{csp})$ and $\cP = \cNP$ over binary words. Then
  for every setting $s$, $\mathrm{frag}(s)$ has the dichotomy and has no intermediate problems.
  \own{}
  \leanfile{Sagitta/Map/Borde/TresInterruptores.lean}
\end{theorem}

\begin{proof}
  \leanok
  The fragment is contained in $\cNP = \cP$, so all its languages are in $\cP$; an intermediate
  problem would give $\cP \neq \cNP$.
\end{proof}

\begin{definition}[Hypothesis: what is missing for Schaefer's theorem, without the separation]
  \label{def:structure-schaefer-sin-separar}
  \lean{Sagitta.Map.Borde.TresInterruptores.SchaeferSinSeparar}
  \leanok
  \uses{def:structure-supuestos, def:structure-dicotomia-defs}
  For $\mathrm{cspDe}$ assigning to each Boolean language $\Gamma$ a language over binary words (the
  encoding of the instances is a parameter), $\mathrm{SchaeferSinSeparar}(\mathrm{cspDe})$ is the
  conjunction of the fields of the hypothesis $\mathrm{Supuestos}$ with membership in $\cP$ and
  $\cNP$-completeness of cslib-frontier in place of the opaque predicates, except the separation
  field, and with the added field that $\mathrm{cspDe}(\Gamma) \in \cNP$ for every finite $\Gamma$.
  This entry says that the definition exists, not that it holds.
  \own{}
  \leanfile{Sagitta/Map/Borde/TresInterruptores.lean}
\end{definition}

\begin{theorem}[The finite Boolean constraint languages have the dichotomy, without assuming $\cP \neq \cNP$]
  \label{thm:structure-schaefer-bz}
  \lean{Sagitta.Map.Borde.TresInterruptores.schaefer_da_BZ_booleano, Sagitta.Map.Borde.TresInterruptores.cspBooleanos}
  \leanok
  \uses{def:structure-schaefer-sin-separar, thm:structure-dicotomia}
  Assume $\mathrm{SchaeferSinSeparar}(\mathrm{cspDe})$. Then the class of languages $\mathrm{cspDe}(\Gamma)$
  for finite Boolean languages $\Gamma$ has the dichotomy.
  \own{}
  \leanfile{Sagitta/Map/Borde/TresInterruptores.lean}
\end{theorem}

\begin{proof}
  \leanok
  If $\cP = \cNP$ every such language is in $\cP$. Otherwise the fields give a
  $\mathrm{Supuestos}$ whose separation field follows from $\cP \neq \cNP$, and
  Theorem~\ref{thm:structure-dicotomia} applies.
\end{proof}

\section{Padding: what it conserves and what it costs}
\label{sec:structure-padding}

The statements of this section are about the padded language $\mathrm{relleno}(g, L)$, which
contains the words made of a framing of a word $x \in L$ followed by $g(|x|)$ ones. Where the
statement is about running time, the model is said: the Turing machines of cslib-frontier, or an
additive cost model in which a decider is a pair of an answer function and a cost function.

\begin{definition}[Padded language]
  \label{def:structure-relleno}
  \lean{Sagitta.Map.Borde.RellenoReloj.rellenar, Sagitta.Map.Borde.RellenoReloj.relleno, Sagitta.Map.Borde.RellenoReloj.largoRelleno}
  \leanok
  For $g : \mathbb{N} \to \mathbb{N}$ and a binary word $x$, $\mathrm{rellenar}(g, x)$ is the framing of
  $x$ (the framing of cslib-frontier) followed by $g(|x|)$ copies of the bit $1$; its length is
  $2|x| + 1 + g(|x|)$. The padded language $\mathrm{relleno}(g, L)$ is the set of words $y$ with $y =
  \mathrm{rellenar}(g, x)$ and $x \in L$, where $x$ is the word framed in $y$.
  \own{}
  \leanfile{Sagitta/Map/Borde/RellenoReloj.lean}
\end{definition}

\begin{theorem}[Polynomial padding is a polynomial equivalence]
  \label{thm:structure-relleno-equivalente}
  \lean{Sagitta.Map.Borde.RellenoReloj.relleno_equivalente}
  \leanok
  \uses{def:structure-relleno, def:structure-dicotomia-defs}
  Let $g$ be computable in polynomial time with input and output in unary, and let $L$ be a
  language with some word $x_0 \notin L$. Then $L$ and $\mathrm{relleno}(g, L)$ are equivalent.
  \own{}
  \leanfile{Sagitta/Map/Borde/RellenoReloj.lean}
\end{theorem}

\begin{proof}
  \leanok
  The map $x \mapsto \mathrm{rellenar}(g, x)$ is a reduction of $L$ to its padded language. The
  reduction back sends a word of the form $\mathrm{rellenar}(g, \mathrm{unframe}\,y)$ to the framed word
  and every other word to the fixed non-member $x_0$.
\end{proof}

\begin{theorem}[Polynomial padding preserves being intermediate]
  \label{thm:structure-relleno-intermedio}
  \lean{Sagitta.Map.Borde.RellenoReloj.relleno_intermedio}
  \leanok
  \uses{def:structure-relleno, thm:structure-relleno-equivalente, thm:structure-intermedio-equiv}
  Let $g$ be computable in polynomial time, in unary, and let $L$ be intermediate. Then
  $\mathrm{relleno}(g, L)$ is intermediate.
  \own{}
  \leanfile{Sagitta/Map/Borde/RellenoReloj.lean}
\end{theorem}

\begin{proof}
  \leanok
  An intermediate language is not the set of all words (that set is in $\cP$), so it has a
  non-member, and Theorem~\ref{thm:structure-relleno-equivalente} and
  Theorem~\ref{thm:structure-intermedio-equiv} apply.
\end{proof}

\begin{definition}[Deciders with a cost, and the padded decider]
  \label{def:structure-decisor}
  \lean{Sagitta.Map.Borde.RellenoReloj.Decisor, Sagitta.Map.Borde.RellenoReloj.Decisor.DecideEn, Sagitta.Map.Borde.RellenoReloj.Decisor.rellenado, Sagitta.Map.Borde.RellenoReloj.inversa}
  \leanok
  \uses{def:structure-relleno}
  A decider is a pair of functions on binary words: the answer and the cost. It decides $L$ in time $T$
  when the answer is $\mathrm{true}$ exactly on the words of $L$ and the cost on $x$ is at most
  $T(|x|)$. For a decider $D$ and functions $g, p$, the padded decider answers $\mathrm{true}$ on $y$
  when $y = \mathrm{rellenar}(g, \mathrm{unframe}\,y)$ and $D$ accepts $\mathrm{unframe}\,y$; its cost on
  $y$ is $p(|y|)$ plus, if $y$ has that form, the cost of $D$ on $\mathrm{unframe}\,y$. The function
  $\mathrm{inversa}(g, m)$ is the greatest $n \le m$ with $2n + 1 + g(n) \le m$ ($0$ if there is
  none). This is an additive cost model, not a machine.
  \own{}
  \leanfile{Sagitta/Map/Borde/RellenoReloj.lean}
\end{definition}

\begin{theorem}[The work of a padded word is the old work plus the guard]
  \label{thm:structure-trabajo}
  \lean{Sagitta.Map.Borde.RellenoReloj.trabajo_no_se_destruye}
  \leanok
  \uses{def:structure-decisor}
  For a decider $D$, functions $g, p$ and a binary word $x$, the cost of the padded decider on
  $\mathrm{rellenar}(g, x)$ is $p(2|x| + 1 + g(|x|)) + \mathrm{cost}_D(x)$.
  \own{}
  \leanfile{Sagitta/Map/Borde/RellenoReloj.lean}
\end{theorem}

\begin{proof}
  \leanok
  Unfold the cost of the padded decider; the word has the right form and $\mathrm{unframe}$ returns
  $x$.
\end{proof}

\begin{theorem}[Translation by exponential padding, in the additive cost model]
  \label{thm:structure-traslacion-exp}
  \lean{Sagitta.Map.Borde.RellenoReloj.traslacion_exponencial}
  \leanok
  \uses{def:structure-decisor}
  Let $D$ be a decider that decides $L$ in time $n \mapsto 2^n$ and let $p : \mathbb{N} \to
  \mathbb{N}$. Then the padded decider, for $g(n) = 2^n$, decides $\mathrm{relleno}(2^n, L)$ in time
  $m \mapsto p(m) + (m + 1)$.
  \own{}
  \leanfile{Sagitta/Map/Borde/RellenoReloj.lean}
\end{theorem}

\begin{proof}
  \leanok
  The general translation gives time $p(m) + T(\mathrm{inversa}(g, m))$ for a monotone $T$, and with $T(n) =
  2^n$ and $g(n) = 2^n$ the second term is at most $m + 1$.
\end{proof}

\begin{definition}[Diluted language and uniform witnesses]
  \label{def:structure-diluido}
  \lean{Sagitta.Map.Borde.RellenoReloj.diluido, Sagitta.Map.Borde.RellenoReloj.TestigosUniformes}
  \leanok
  For a language $A$ and a set $S \subseteq \mathbb{N}$ of lengths, $\mathrm{diluido}(A, S)$ is the set of
  words of $A$ whose length is in $S$. The language $A$ has \emph{uniform witnesses} from $N$ when
  some function $w$, computable in polynomial time with input in unary and output a binary word, has
  for every $n \ge N$ a word $w(n)$ of length $n$ in $A$.
  \own{}
  \leanfile{Sagitta/Map/Borde/RellenoReloj.lean}
\end{definition}

\begin{theorem}[Deciding a diluted language in $\cP$ needs a clock for the set of lengths]
  \label{thm:structure-reloj-diluido}
  \lean{Sagitta.Map.Borde.RellenoReloj.reloj_de_diluido}
  \leanok
  \uses{def:structure-diluido}
  Let $A$ have uniform witnesses from $N$ and let $S \subseteq \mathbb{N}$. If $\mathrm{diluido}(A, S)
  \in \cP$ over binary words, then some $c : \mathbb{N} \to \mathrm{Bool}$ computable in polynomial time
  with input in unary satisfies $c(n) = \mathrm{true} \iff n \in S$ for every $n \ge N$.
  \own{}
  \leanfile{Sagitta/Map/Borde/RellenoReloj.lean}
\end{theorem}

\begin{proof}
  \leanok
  Take $c(n)$ to be the decider of $\mathrm{diluido}(A, S)$ applied to $w(n)$; since $w(n) \in A$ and
  $|w(n)| = n$, the answer is whether $n \in S$.
\end{proof}

\begin{definition}[Strict reduction, quantized and dense classes]
  \label{def:structure-denso}
  \lean{Sagitta.Map.Borde.RellenoReloj.Estricta, Sagitta.Map.Borde.RellenoReloj.Cuantizado, Sagitta.Map.Borde.RellenoReloj.Denso}
  \leanok
  \uses{def:structure-dicotomia-defs}
  $A <_{\mathrm{p}} B$ (\texttt{Estricta}) means $A \red B$ and not $B \red A$. A class $F$ is
  \emph{quantized} when it has the dichotomy. It is \emph{dense} when for all $A, B \in F$ with $B
  \notin \cP$ and $A <_{\mathrm{p}} B$ there is $C \in F$ with $A <_{\mathrm{p}} C <_{\mathrm{p}} B$.
  \own{}
  \leanfile{Sagitta/Map/Borde/RellenoReloj.lean}
\end{definition}

\begin{theorem}[$\cNP$ is dense]
  \label{thm:structure-denso-np}
  \lean{Sagitta.Map.Borde.RellenoReloj.denso_NP}
  \leanok
  \uses{def:structure-denso, def:NP}
  The class $\cNP$ of languages over binary words is dense.
  \source{\cite{Ladner75}.}
  \fidelity{Not checked against the printed source.}
  \leanfile{Sagitta/Map/Borde/RellenoReloj.lean}
\end{theorem}

\begin{proof}
  \leanok
  It is the density theorem of cslib-frontier ($A <_{\mathrm{p}} B$ with $B \notin \cP$ in $\cNP$ has a
  language of $\cNP$ strictly between them), stated with the two strict reductions.
\end{proof}

\begin{theorem}[A dense class with a non-trivial member of $\cP$ and a complete problem is not quantized]
  \label{thm:structure-denso-no-cuantizado}
  \lean{Sagitta.Map.Borde.RellenoReloj.denso_no_cuantizado}
  \leanok
  \uses{def:structure-denso}
  Assume $\cP \neq \cNP$ over binary words and let $F$ be a dense class. Let $A \in F \cap \cP$ with
  words $a_1 \in A$ and $a_0 \notin A$, and let $B \in F$ be $\cNP$-complete. Then $F$ is not
  quantized.
  \own{}
  \leanfile{Sagitta/Map/Borde/RellenoReloj.lean}
\end{theorem}

\begin{proof}
  \leanok
  $B \notin \cP$ because $\cP \neq \cNP$, and $A <_{\mathrm{p}} B$; density gives a $C \in F$ strictly
  between, which is neither in $\cP$ nor complete, so the dichotomy fails.
\end{proof}

\begin{definition}[A class with a clock]
  \label{def:structure-admite-reloj}
  \lean{Sagitta.Map.Borde.RellenoReloj.AdmiteReloj}
  \leanok
  \uses{def:structure-dicotomia-defs}
  A class $F$ \emph{admits a clock} when for every $B \in F$ and every $G \in \cP$ there is $C \in F$
  equivalent to the set of words of $B$ that are not in $G$.
  \own{}
  \leanfile{Sagitta/Map/Borde/RellenoReloj.lean}
\end{definition}

\begin{theorem}[A class with a clock and a complete problem has intermediate problems, if $\cP \neq \cNP$]
  \label{thm:structure-medio-reloj}
  \lean{Sagitta.Map.Borde.RellenoReloj.medio_de_reloj}
  \leanok
  \uses{def:structure-admite-reloj, thm:structure-ladner-palabras}
  Assume $\cP \neq \cNP$ over binary words. Let $F$ be a class that admits a clock and let $B \in F$ be
  $\cNP$-complete. Then $F$ has intermediate problems.
  \own{}
  \leanfile{Sagitta/Map/Borde/RellenoReloj.lean}
\end{theorem}

\begin{proof}
  \leanok
  The set $B$ is not in $\cP$. The downward form of Ladner's theorem of cslib-frontier gives $G \in
  \cP$ with the words of $B$ outside $G$ in $\cNP$, outside $\cP$ and strictly below $B$; the clock
  puts an equivalent language in $F$, and it is not complete because a reduction from $B$ would
  contradict strictness.
\end{proof}

\begin{definition}[Hypothesis: switching off an interrupter gives a clock]
  \label{def:structure-interruptor-reloj}
  \lean{Sagitta.Map.Borde.RellenoReloj.InterruptorEsReloj}
  \leanok
  \uses{def:structure-admite-reloj, def:structure-interruptores}
  For an assignment $\mathrm{frag}$ of a class to each setting of the switches,
  $\mathrm{InterruptorEsReloj}(\mathrm{frag})$ says that $\mathrm{frag}(s)$ admits a clock for every
  setting $s \neq \mathrm{encendido}$. This entry says that the definition exists, not that it holds.
  \own{}
  \leanfile{Sagitta/Map/Borde/RellenoReloj.lean}
\end{definition}

\begin{theorem}[From the captures of Feder and Vardi, switching off gives a clock]
  \label{thm:structure-reloj-de-fv}
  \lean{Sagitta.Map.Borde.RellenoReloj.interruptorEsReloj_de_FV}
  \leanok
  \uses{def:structure-feder-vardi, def:structure-interruptor-reloj, thm:structure-captura-apagado}
  Assume $\mathrm{FederVardi}(\mathrm{frag}, \mathrm{csp})$. Then $\mathrm{InterruptorEsReloj}(\mathrm{frag})$.
  \own{}
  \leanfile{Sagitta/Map/Borde/RellenoReloj.lean}
\end{theorem}

\begin{proof}
  \leanok
  A class contained in $\cNP$ that captures $\cNP$ admits a clock: the set of words of $B$ outside $G$ is
  in $\cNP$, and the capture gives an equivalent member of the class. The fragments with a switch off
  capture $\cNP$ by Theorem~\ref{thm:structure-captura-apagado}.
\end{proof}

\begin{theorem}[With all switches on there is no clock, under the assumptions]
  \label{thm:structure-encendido-sin-reloj}
  \lean{Sagitta.Map.Borde.RellenoReloj.encendido_sin_reloj}
  \leanok
  \uses{def:structure-feder-vardi, def:structure-admite-reloj, thm:structure-medio-reloj, thm:structure-encendido-dicotomia}
  Assume $\mathrm{FederVardi}(\mathrm{frag}, \mathrm{csp})$, that $\mathrm{csp}$ has the dichotomy, and that
  $\cP \neq \cNP$ over binary words. If $B \in \mathrm{frag}(\mathrm{encendido})$ is $\cNP$-complete, then
  $\mathrm{frag}(\mathrm{encendido})$ does not admit a clock.
  \own{}
  \leanfile{Sagitta/Map/Borde/RellenoReloj.lean}
\end{theorem}

\begin{proof}
  \leanok
  A clock would give intermediate problems in the fragment (Theorem~\ref{thm:structure-medio-reloj}),
  contradicting its dichotomy (Theorem~\ref{thm:structure-encendido-dicotomia}).
\end{proof}

\begin{definition}[Exponential padding of the diagonal language, and the first law as a statement]
  \label{def:structure-relleno-exp}
  \lean{Sagitta.Map.Borde.RellenoExponencial.expo, Sagitta.Map.Borde.RellenoExponencial.diagonalL, Sagitta.Map.Borde.RellenoExponencial.diagonalRellena, Sagitta.Map.Borde.RellenoExponencial.PrimeraLey}
  \leanok
  \uses{def:structure-relleno}
  $\mathrm{expo}(n) = 2^n$. $\mathrm{diagonalL}$ is the set of words $x$ with $\mathrm{diagonal}(x) =
  \mathrm{true}$, the diagonal language of the time hierarchy of cslib-frontier;
  $\mathrm{diagonalRellena}$ is $\mathrm{relleno}(\mathrm{expo}, \mathrm{diagonalL})$. For $g : \mathbb{N}
  \to \mathbb{N}$, $\mathrm{PrimeraLey}(g)$ says that for every language $L$ with some non-member,
  $L \in \cP$ if and only if $\mathrm{relleno}(g, L) \in \cP$, over binary words.
  \own{}
  \leanfile{Sagitta/Map/Borde/RellenoExponencial.lean}
\end{definition}

\begin{theorem}[A language outside $\cP$ whose exponentially padded version is in $\cP$]
  \label{thm:structure-dificil-relleno-facil}
  \lean{Sagitta.Map.Borde.RellenoExponencial.existe_dificil_con_relleno_facil}
  \leanok
  \uses{def:structure-relleno-exp, def:P}
  There is a language $A$ of binary words with $A \notin \cP$ and $\mathrm{relleno}(\mathrm{expo}, A) \in
  \cP$, for the Turing machines of cslib-frontier.
  \own{}
  \leanfile{Sagitta/Map/Borde/RellenoExponencial.lean}
\end{theorem}

\begin{proof}
  \leanok
  Take $A = \mathrm{diagonalL}$, which is not in $\cP$ by the time hierarchy of cslib-frontier. A machine
  first checks cheaply that the framed word is shorter than $\log_2$ of the input, then that the input
  is the padded word, and then runs the diagonal on the framed word in time at most polynomial in the
  length of the input.
\end{proof}

\begin{theorem}[The first law holds for polynomial padding and fails for exponential padding]
  \label{thm:structure-primera-ley}
  \lean{Sagitta.Map.Borde.RellenoExponencial.primeraLey_de_polinomial, Sagitta.Map.Borde.RellenoExponencial.no_primeraLey_expo}
  \leanok
  \uses{def:structure-relleno-exp, thm:structure-relleno-equivalente, thm:structure-dificil-relleno-facil}
  If $g$ is computable in polynomial time, with input and output in unary, then $\mathrm{PrimeraLey}(g)$
  holds. The statement $\mathrm{PrimeraLey}(\mathrm{expo})$ does not hold.
  \own{}
  \leanfile{Sagitta/Map/Borde/RellenoExponencial.lean}
\end{theorem}

\begin{proof}
  \leanok
  The first statement is the equivalence of Theorem~\ref{thm:structure-relleno-equivalente}. For the
  second, the diagonal language is not in $\cP$ and has a non-member, while its exponential padding is
  in $\cP$ by Theorem~\ref{thm:structure-dificil-relleno-facil}.
\end{proof}

\section{Universal enumerations, the diagonal and the lazy diagonal}
\label{sec:structure-diagonal}

The statements of this section are about enumerations of sets or functions, in a model with no
machine; the last two are about the class $\cP$ of cslib-frontier.

\begin{definition}[Universal family, diagonal of an enumeration]
  \label{def:structure-universal}
  \lean{Sagitta.Map.Borde.UniversalDiagonal.Universal, Sagitta.Map.Borde.UniversalDiagonal.diagonal}
  \leanok
  A family $e : A \to A \to B$ is \emph{universal} when it is surjective as a map from $A$ to the
  functions $A \to B$: every function $A \to B$ is a row of $e$. For an enumeration $M : \mathbb{N}
  \to \mathrm{Set}\,\alpha$ and a map $\mathrm{idx} : \alpha \to \mathbb{N}$, the diagonal is the set of
  $a \in \alpha$ with $a \notin M(\mathrm{idx}(a))$.
  \own{}
  \leanfile{Sagitta/Map/Borde/UniversalDiagonal.lean}
\end{definition}

\begin{theorem}[No universal family if some self-map of the codomain has no fixed point]
  \label{thm:structure-sin-universal}
  \lean{Sagitta.Map.Borde.UniversalDiagonal.sin_enumeracion_universal}
  \leanok
  \uses{def:structure-universal}
  Let $f : B \to B$ with $f(b) \neq b$ for every $b$. Then no family $e : A \to A \to B$ is universal.
  \source{\cite{Lawvere69}.}
  \fidelity{Not checked against the printed source.}
  \leanfile{Sagitta/Map/Borde/UniversalDiagonal.lean}
\end{theorem}

\begin{proof}
  \leanok
  A universal family forces every self-map of $B$ to have a fixed point (Lawvere's fixed-point
  theorem, which is that of Mathlib), contradicting the hypothesis.
\end{proof}

\begin{theorem}[The diagonal of an enumeration is not in the enumeration]
  \label{thm:structure-diagonal-fuera}
  \lean{Sagitta.Map.Borde.UniversalDiagonal.diagonal_fuera}
  \leanok
  \uses{def:structure-universal}
  Let $M : \mathbb{N} \to \mathrm{Set}\,\alpha$ and let $\mathrm{idx} : \alpha \to \mathbb{N}$ be
  surjective. Then the diagonal of $M$ is not in the range of $M$.
  \own{}
  \leanfile{Sagitta/Map/Borde/UniversalDiagonal.lean}
\end{theorem}

\begin{proof}
  \leanok
  The diagonal is in the range exactly when the family $(a, b) \mapsto b \in M(\mathrm{idx}(a))$
  represents the function that negates its own diagonal, and Lawvere's argument with negation, which
  has no fixed point, rules this out.
\end{proof}

\begin{theorem}[The diagonal of the presentation of $\cP$ is not in $\cP$]
  \label{thm:structure-diagonal-fuera-p}
  \lean{Sagitta.Map.Borde.UniversalDiagonal.diagonal_fuera_de_P}
  \leanok
  \uses{thm:structure-diagonal-fuera, def:P}
  There is a function $F : \mathbb{N} \times \mathrm{List}\,\mathrm{Bool} \to \mathrm{Bool}$ such that the set
  of binary words $w$ with $F(\mathrm{index}(w), w) = \mathrm{false}$ is not in $\cP$, where
  $\mathrm{index}$ is the bijection between binary words and natural numbers of cslib-frontier.
  \own{}
  \leanfile{Sagitta/Map/Borde/UniversalDiagonal.lean}
\end{theorem}

\begin{proof}
  \leanok
  The class $\cP$ is recursively presentable in cslib-frontier, by an $F$ whose rows are the sets of
  the class; its diagonal is outside the class by Theorem~\ref{thm:structure-diagonal-fuera}.
\end{proof}

\begin{theorem}[Lazy diagonal: a mixed set outside two enumerated classes]
  \label{thm:structure-schoning-abstracto}
  \lean{Sagitta.Map.Borde.UniversalDiagonal.schoning_abstracto}
  \leanok
  Let $M_1, M_2 : \mathbb{N} \to \mathrm{Set}\,\mathbb{N}$ and let $C_1 = \mathrm{range}\,M_1$ and $C_2 =
  \mathrm{range}\,M_2$. Assume that $C_1$ and $C_2$ are closed under finite variation, and let $A_1,
  A_2 \subseteq \mathbb{N}$ with $A_1 \notin C_1$ and $A_2 \notin C_2$. Then there is $G \subseteq
  \mathbb{N}$ such that the set equal to $A_1$ on $G$ and to $A_2$ outside $G$ is neither in $C_1$ nor
  in $C_2$.
  \source{\cite{Schoning82}.}
  \fidelity{Not checked against the printed source.}
  \leanfile{Sagitta/Map/Borde/UniversalDiagonal.lean}
\end{theorem}

\begin{proof}
  \leanok
  Closure under finite variation makes each difference $A_1 \mathbin{\triangle} M_1(i)$ and $A_2
  \mathbin{\triangle} M_2(i)$ infinite; choose a strictly increasing sequence of witnesses, alternating
  between the two enumerations, and let $G$ be the set of the even-indexed witnesses.
\end{proof}

\begin{definition}[Hypothesis: a class is recursively presentable]
  \label{def:structure-enumeracion}
  \lean{Sagitta.Map.Borde.UniversalDiagonal.Enumeracion}
  \leanok
  For a class $C$ of sets of binary words, $\mathrm{Enumeracion}(C)$ says that $C$ is recursively
  presentable for the binary encoding of words, in the sense of cslib-frontier. This entry says that
  the definition exists, not that it holds for a given $C$.
  \own{}
  \leanfile{Sagitta/Map/Borde/UniversalDiagonal.lean}
\end{definition}

\begin{theorem}[Ladner's theorem downward, from two enumerations]
  \label{thm:structure-ladner-enumeraciones}
  \lean{Sagitta.Map.Borde.UniversalDiagonal.ladner_de_enumeraciones}
  \leanok
  \uses{def:structure-enumeracion}
  Let $B$ be a decidable set of binary words with $B \notin \cP$. Assume $\mathrm{Enumeracion}$ of the
  polynomial-time degree of $B$ (the set of languages equivalent to $B$ under reductions in both
  directions) and of the class $\cP$. Then there is $G \in \cP$ such that the set of words of $B$
  outside $G$ is not in $\cP$, reduces to $B$ and is not reducible from $B$.
  \own{}
  \leanfile{Sagitta/Map/Borde/UniversalDiagonal.lean}
\end{theorem}

\begin{proof}
  \leanok
  It applies the effective lazy diagonal of cslib-frontier to the two enumerations, using the closure of
  both classes under finite variation and that the empty set is not in the degree of $B$.
\end{proof}

\begin{theorem}[Ladner's theorem downward, without hypotheses]
  \label{thm:structure-ladner-universal}
  \lean{Sagitta.Map.Borde.UniversalDiagonal.ladner_cslib_universal}
  \leanok
  \uses{thm:structure-ladner-enumeraciones, thm:ladner}
  Let $B$ be a decidable set of binary words with $B \notin \cP$. Then there is $G \in \cP$ such that the
  set of words of $B$ outside $G$ is not in $\cP$, reduces to $B$ and is not reducible from $B$.
  \source{\cite{Ladner75,Schoning82}.}
  \fidelity{Not checked against the printed source.}
  \leanfile{Sagitta/Map/Borde/UniversalDiagonal.lean}
\end{theorem}

\begin{proof}
  \leanok
  Theorem~\ref{thm:structure-ladner-enumeraciones} with the two enumerations that cslib-frontier proves
  to exist (those of $\cP$ and of the degree of $B$).
\end{proof}

\begin{proposition}[Without a universal family there is no diagonal: the constants]
  \label{prop:structure-constantes}
  \lean{Sagitta.Map.Borde.UniversalDiagonal.constantes, Sagitta.Map.Borde.UniversalDiagonal.constantes_no_universal}
  \leanok
  \uses{def:structure-universal}
  Let $\mathrm{constantes} : \mathrm{Bool} \to \mathrm{Bool} \to \mathrm{Bool}$ be the family $a \mapsto (b
  \mapsto a)$. Then $\mathrm{constantes}$ is not universal.
  \own{}
  \leanfile{Sagitta/Map/Borde/UniversalDiagonal.lean}
\end{proposition}

\begin{proof}
  \leanok
  If the identity were a row $a$, evaluating at $\mathrm{false}$ and at $\mathrm{true}$ would give
  $a = \mathrm{false}$ and $a = \mathrm{true}$.
\end{proof}
